\section{Fibre constancy}\label{sec:fc}

This chapter proves Theorem~\ref{thm:gl-fibre-constancy}, the claim ``$\Phi$ is constant on every fibre of $\beta^\an$''
in Step~1 of \cite[Theorem~6.1]{FHX}. As in \cite{FHX}, the fibre is the analytification of a proper connected
$\mathcal H(t)$-scheme, and the restriction of $\Phi$ is plurisubharmonic in the sense of curves \cite[Definition~4.1]{GR25}.
We use two deviations from \cite{GR25}. Subharmonicity on curves is tested through local norms to the Berkovich
line (Definition~\ref{def:fc-subharmonic}), so potential theory is only needed on the line. The maximum principle on
the fibre is proved by chains of algebraic curves through rigid points (Lemma~\ref{lem:fc-curve-chains}), instead of
the local maximum principle \cite[Theorem~6.1]{GR25} and Berkovich's connectedness theorem \cite[3.4.8]{Berkovich90}.
See Remark~\ref{rmk:fc-gr25}.

\subsection{Setting and reductions}\label{subsec:fc-setting}

\begin{convention}\label{conv:fc-setting}
\notready
\uses{conv:gl-setting, def:pg-blowup, def:pg-blowup-psh, def:convex-pl}
$k$ is countable, algebraically closed and trivially valued; $T'=V(I)\subseteq\PP^N$ is projectively normal, $B$ is a
blowup of $T'$ as in Definition~\ref{def:pg-blowup}, $\mu\colon B^\an\to T'^\an$, $a\in\Q_{>0}$, $\Phi\in\PSH(B,(a,0))$.
A \emph{complete field} is a complete non-Archimedean field $F$ with a $k$-algebra structure such that $|c|=1$ for
$c\in k^\times$. For an $F$-scheme $X$ of finite type, $X^\an$ is the space of pairs (scheme point, real valuation of
its residue field extending that of $F$); for a complete extension $F'/F$, $\pi_{F'/F}\colon X_{F'}^\an\to X^\an$ is
restriction. A \emph{curve} over $F$ is a separated $F$-scheme of finite type of pure dimension one. Locally psh
functions on open subsets of $X^\an$ are those of Definition~\ref{def:convex-pl}, with units on Zariski open subsets.
\end{convention}

\begin{lemma}[Affine psh pullback under scalar extension]\label{lem:fc-affine-base-change-psh}
\leanok
\lean{EnvelopeProperty.Berkovich.RelativeSpectrum.IsConvexPLOn.fc_comp_baseChangeRestriction, EnvelopeProperty.Berkovich.RelativeSpectrum.IsDecrConvexPLLimitOn.fc_comp_baseChangeRestriction, EnvelopeProperty.Berkovich.RelativeSpectrum.IsLocallyPSHOn.fc_comp_baseChangeRestriction}
\uses{conv:fc-setting, lem:local-basic}
Let $R\to S$ be a homomorphism of commutative rings with compatible real-valued valuations $v,w$, and $A$ an
$R$-algebra. Pullback along the restriction map
$\pi\colon\operatorname{Berk}_w(S\otimes_RA)\to\operatorname{Berk}_v(A)$ preserves convex PL functions,
decreasing limits of convex PL functions, and locally psh functions on open subsets.
\end{lemma}

\begin{proof}
Replace each coefficient $a_j\in A$ in a finite convex PL presentation by $1\otimes a_j$.
Since $|a_j|_{\pi(z)}=|1\otimes a_j|_z$, this preserves the rational affine logarithmic expressions,
their maxima, and their nonvanishing loci. Pullback commutes with pointwise infima and preserves decreasing
sequences. Finally, $\pi$ is continuous, so open neighborhoods in a local psh presentation pull back to open
neighborhoods.
\end{proof}

\begin{lemma}[Rational scalar lifts]\label{lem:fc-rational-scalar-lift}
\leanok
\lean{EnvelopeProperty.Berkovich.RelativeSpectrum.fcRationalLiftEval, EnvelopeProperty.Berkovich.RelativeSpectrum.fcRationalLiftEval_tmul, EnvelopeProperty.Berkovich.RelativeSpectrum.fcRationalLiftEval_surjective, EnvelopeProperty.Berkovich.RelativeSpectrum.fcRationalLiftQuotientEquiv, EnvelopeProperty.Berkovich.RelativeSpectrum.fcRationalLift, EnvelopeProperty.Berkovich.RelativeSpectrum.fcRationalLift_baseChangeRestriction}
\uses{conv:fc-setting}
Let $F\to L$ be a homomorphism of commutative rings with compatible real-valued valuations $v,w$.
For every $F$-algebra $A$ and $F$-algebra homomorphism $e\colon A\to F$, the map
$e_L\colon L\otimes_FA\to L$, $c\otimes a\mapsto c\,e(a)$, is a surjective $L$-algebra homomorphism.
Its valuation $w\circ e_L$ restricts to the original rational seminorm $v\circ e$ on $A$.
\end{lemma}

\begin{proof}
The two maps $L\to L$ and $A\xrightarrow{e}F\to L$ agree on $F$ and have commuting images,
so induce the stated tensor homomorphism. Its value on $c\otimes1$ is $c$, proving surjectivity.
Its value on $1\otimes a$ is $e(a)$ viewed in $L$, and compatibility of $v,w$ proves the last assertion.
\end{proof}

\begin{lemma}[Lifting points through completion]\label{lem:fc-completion-lift}
\leanok
\lean{EnvelopeProperty.Berkovich.RelativeSpectrum.fcCompletionCoefficientHom, EnvelopeProperty.Berkovich.RelativeSpectrum.fcCompletionCoefficientHom_coe, EnvelopeProperty.Berkovich.RelativeSpectrum.continuous_fcCompletionCoefficientHom, EnvelopeProperty.Berkovich.RelativeSpectrum.norm_fcCompletionCoefficientHom, EnvelopeProperty.Berkovich.RelativeSpectrum.fcCompletionEval, EnvelopeProperty.Berkovich.RelativeSpectrum.fcCompletionEval_tmul, EnvelopeProperty.Berkovich.RelativeSpectrum.fcCompletion_val_algebraMap, EnvelopeProperty.Berkovich.RelativeSpectrum.fcCompletionLift, EnvelopeProperty.Berkovich.RelativeSpectrum.fcCompletionLift_baseChangeRestriction, EnvelopeProperty.Berkovich.RelativeSpectrum.fc_baseChangeRestriction_completion_surjective}
\uses{conv:fc-setting}
Let $F$ be a non-Archimedean normed field and $\widehat F$ its completion. For every commutative $F$-algebra $A$,
the restriction map from seminorms on $\widehat F\otimes_FA$ prescribing the norm of $\widehat F$ to seminorms on
$A$ prescribing the norm of $F$ is surjective.
\end{lemma}

\begin{proof}
For a given point $x$, the coefficient map $F\to\mathcal H(x)$ is isometric. It therefore extends to a continuous
homomorphism $\widehat F\to\mathcal H(x)$, which remains isometric by density. Combining this map with the original
evaluation $A\to\mathcal H(x)$ gives a homomorphism $\widehat F\otimes_FA\to\mathcal H(x)$. Pulling back the norm
of $\mathcal H(x)$ gives the required seminorm: on the left tensor factor it equals the whole norm of
$\widehat F$, while on $1\otimes a$ it equals $|a|_x$.
\end{proof}

\begin{lemma}[Completed algebraic scalar extensions]\label{lem:fc-algebraic-completion-lift}
\leanok
\lean{EnvelopeProperty.Berkovich.RelativeSpectrum.fcCompletedExtension_val_algebraMap, EnvelopeProperty.Berkovich.RelativeSpectrum.fc_baseChangeRestriction_algebraicCompletion_surjective}
\uses{lem:fc-completion-lift}
Let $K$ be complete non-Archimedean, and $L/K$ an algebraic normed extension. For every commutative $K$-algebra $A$,
scalar restriction from the norm-prescribed spectrum of $\widehat L\otimes_KA$ onto that of $A$ over $K$ is surjective.
\end{lemma}

\begin{proof}
An algebraic field extension makes $A\to L\otimes_KA$ integral and injective. Integral point extension, together
with uniqueness of extending the norm of the complete field $K$, supplies a point over the original seminorm and
prescribing the norm of $L$. Lemma~\ref{lem:fc-completion-lift} lifts this point to
$\widehat L\otimes_L(L\otimes_KA)$. The actual scalar tensor cancellation isomorphism with
$\widehat L\otimes_KA$ retains every original $A$ value and every completed coefficient value, giving the claimed lift.
\end{proof}

\begin{lemma}[Extension of scalars]\label{lem:fc-base-change}
\leanok
\lean{EnvelopeProperty.Berkovich.SchemeScalarBaseChange.continuous_nativeMap, EnvelopeProperty.Berkovich.SchemeScalarBaseChange.isClosedMap_nativeMap, EnvelopeProperty.Berkovich.SchemeScalarBaseChange.isLocallyPSHOn_comp_nativeMap, EnvelopeProperty.Berkovich.SchemeScalarBaseChange.nativeMap_surjective_algebraicCompletion, EnvelopeProperty.Berkovich.SchemeScalarBaseChange.nativeMap_surjective_gauss, EnvelopeProperty.Berkovich.SchemeScalarBaseChange.fcRationalScalarLift_nativeMap, EnvelopeProperty.Berkovich.SchemeScalarBaseChange.fcScalarChangeRationalSection_nativeProjection}
\uses{conv:fc-setting, lem:fc-scheme-scalar-pullback, lem:fc-scheme-scalar-proper, lem:fc-scheme-scalar-algebraic-completion-lift, lem:fc-scheme-gauss-scalar-lift, lem:fc-scheme-rational-point-lift, lem:fc-scheme-rational-section-base-change, lem:fc-scheme-nontrivial-complete-scalar-lift}
Let $X$ be an $F$-scheme of finite type and $F'/F$ a complete extension.
\begin{enumerate}
\item $\pi_{F'/F}$ is continuous and closed, and the pullback of a locally psh function is locally psh.
\item $\pi_{F'/F}$ is surjective if $F'$ is the completion of an algebraic extension of $F$, or the field
$\mathcal H(\eta_r)$ of the Gauss point of radius $r>0$ of $\A^1_F$.
\item For $x\in X^\an$ let $F_x$ be the completed algebraic closure of $\mathcal H(x)$. There is an $F_x$-rational point
$\tilde x$ of $X_{F_x}$ with $\pi_{F_x/F}(\tilde x)=x$. Every $F$-rational point $c$ of $X$ gives an $F'$-rational point
$\tilde c$ over $c$.
\item There is a complete, algebraically closed, nontrivially valued extension $F_1/F$ such that $\pi_{F_1/F}$ is
surjective.
\end{enumerate}
\end{lemma}

\begin{proof}
(1) On actual affine charts the whole scalar restriction map is the genuine
tensor-spectrum map. Ultrafilter lifting proves properness there; actual open
chart pullbacks give properness of the whole scheme map. Thus it is continuous
and closed. Regular units and their whole logarithmic values pull back along
the genuine scheme projection, preserving local convex PL presentations and
their decreasing limits.

(2) For an algebraic extension, extend the actual point through its integral
coefficient algebra and then its completion. The whole projection and all
original coefficients are preserved. For a Gauss field, the actual weighted
Gauss tensor seminorm extends every prescribed original point; pass through
its genuine completed field and actual affine chart. This constructs a whole
lift and proves surjectivity.

(3) The canonical completed stalk residue-field point gives the actual base
coefficient action on $\mathcal H(x)$. Its completed algebraic closure preserves
every norm. The genuine scheme pullback universal property constructs an actual
$F_x$-rational section with whole original projection $x$. The same universal
property lifts each original $F$-rational section to any normed extension.

(4) Take the actual completed algebraic closure of the radius-$1/2$ Gauss field.
The inverse Gauss variable has norm $2$, proving nontriviality. The actual
scalar-tower scheme morphism identifies the entire iterated projection with
the single final scalar projection. Compose the two constructed surjections
from (2), preserving the whole original point.
\end{proof}

\begin{lemma}[The completed-field affine fibre]\label{lem:fc-affine-fibre-model}
\leanok
\lean{EnvelopeProperty.Berk.fcAbsoluteFibreHomeomorph, EnvelopeProperty.Berk.fcCompletedFibreHomeomorph, EnvelopeProperty.Berk.fcCompletedFibreHomeomorph_val, EnvelopeProperty.Berk.fcCompletedFibreHomeomorph_val_tmul, EnvelopeProperty.Berk.fcCompletedFibreInclusion, EnvelopeProperty.Berk.continuous_fcCompletedFibreInclusion, EnvelopeProperty.Berk.fcNativeCompletedFibreHomeomorph, EnvelopeProperty.Berk.fcNativeCompletedFibreHomeomorph_projection}
\uses{conv:fc-setting}
Let $A\to C$ be a homomorphism of $k$-algebras, and $t\in(\Spec A)^\an$.
The space of seminorms on $\mathcal H(t)\otimes_AC$ restricting to the norm of $\mathcal H(t)$ is homeomorphic to
the fibre over $t$ of $(\Spec C)^\an\to(\Spec A)^\an$. The homeomorphism $j_t$ satisfies
$|c|_{j_t(z)}=|1\otimes c|_z$ and $|h\otimes c|_z=\lVert h\rVert\,|c|_{j_t(z)}$.
\end{lemma}

\begin{proof}
The completed-residue-field tensor fibre homeomorphism for unrestricted seminorm spectra gives these maps and
their topology. A seminorm over the point $t$ is automatically trivial on $k^\times$, because its value on
$c\in k^\times$ equals the value at $t$. Thus restricting to the prescribed trivial coefficient locus changes
neither the fibre nor its topology. The identities follow by evaluating the right tensor inclusion and using
multiplicativity with the prescribed norm on the left tensor factor. The actual affine scheme pullback
$\Spec C\times_{\Spec A}\Spec\mathcal H(t)$ is isomorphic to $\Spec(\mathcal H(t)\otimes_AC)$,
compatibly with both projections. Thus the same homeomorphism identifies the original analytic fibre with the
fixed-norm native points of this scheme-theoretic pullback, retaining the whole original projection.
\end{proof}

\begin{lemma}[Compact blowup fibres]\label{lem:fc-fibre-compact}
\leanok
\lean{EnvelopeProperty.ProjSubvariety.BlowupCenter.fc_isCompact_fibre}
\uses{lem:pg-blowup-basic}
Every fibre of $\mu\colon B^\an\to T'^\an$ is compact.
\end{lemma}

\begin{proof}
The total analytic blowup is compact, the target is Hausdorff, and $\mu$ is continuous. Thus the inverse image of
a point is closed in a compact space.
\end{proof}

\begin{lemma}[Actual Segre polynomial restriction]\label{lem:fc-segre-polynomials}

\uses{lem:pg-bigraded-extension}
Put $D=(M+1)(M'+1)-1$, order the pairs $(i,j)$ as the $D+1$ homogeneous coordinates, and define
$\sigma^*\colon k[z_{ij}]\to k[x_i,y_j]$ by $z_{ij}\mapsto x_i y_j$.
Every form of degree $d$ maps to bidegree $(d,d)$, and
$(\sigma^*F)_{(d,d)}=\sigma^*(F_d)$ for every polynomial $F$.
For every normalized bigraded point $v$, restriction along this actual polynomial homomorphism gives
$a(v)\in(\PP^D)^\an$ with $|F|_{a(v)}=|\sigma^*F|_v$ for every polynomial $F$.
The map $a$ is continuous, and all quadratic Segre minors vanish at $a(v)$.
\end{lemma}

\begin{proof}
A substituted monomial of degree $d$ has $x$-degree and $y$-degree both $d$.
Monomial induction then proves that homogeneous components commute with substitution, so the restricted
seminorm is graded. Each coordinate product has value at most one, and a pair of value-one coordinates gives
a product of value one, proving normalization. Evaluation of each fixed substituted polynomial is continuous.
The quadratic minors pull back to the zero polynomial.
\end{proof}

\begin{lemma}[Recovery from Segre coordinates]\label{lem:fc-segre-embedding}
\leanok
\uses{lem:fc-segre-polynomials}
The actual map $a$ of Lemma~\ref{lem:fc-segre-polynomials} is a closed topological embedding.
For a chosen pair $(i,j)$ let $\rho_{ij}\colon k[x,y]\to k[z]$ send $x_l$ to $z_{lj}$ and $y_m$ to $z_{im}$.
For every form $F$ of bidegree $(p,q)$ one has the polynomial identity
$\sigma^*\rho_{ij}(F)=F x_i^q y_j^p$. Thus, when $|x_i|_v=|y_j|_v=1$,
$|\rho_{ij}(F)|_{a(v)}=|F|_v$.
\end{lemma}

\begin{proof}
The polynomial identity follows first for monomials, by induction on their exponent functions, and then for
bihomogeneous forms by weighted homogeneous induction. If $a(v)=a(w)$, choose value-one coordinates $x_i,y_j$
at $v$. Their product has value one at $w$; since both factors have value at most one there, each factor has
value one. The recovery identity gives equality of $v,w$ on every bihomogeneous form, hence on every polynomial.
Thus $a$ is injective. A continuous injection from the compact bigraded point space to the Hausdorff projective
point space is a closed embedding.
\end{proof}

\begin{lemma}[The actual native Segre point map]\label{lem:fc-segre-native}
\leanok
\uses{lem:fc-segre-embedding}
Composing $a$ with the actual native polynomial Proj homeomorphism gives a closed embedding from the bigraded
point space into the native valued points of the genuine scheme $\operatorname{Proj}k[z_{ij}]$, prescribing the
trivial norm of $k$. On a coordinate chart $z_{ij}\ne0$ with $|x_i|_v=|y_j|_v=1$, the actual homogeneous section
$F/z_{ij}^d$ has value $|\sigma^*F|_v$ for every form $F$ of degree $d$.
\end{lemma}

\begin{proof}
The polynomial Proj homeomorphism preserves the original homogeneous chart sections. The pair of value-one
coordinates gives $|z_{ij}|_{a(v)}=1$, so its section-evaluation formula and the exact polynomial-value identity
of Lemma~\ref{lem:fc-segre-polynomials} give the claim. Composition with this homeomorphism preserves the proved
closed embedding. Identification of the image with the actual global blowup scheme is a separate obligation.
\end{proof}

\begin{lemma}[The fibre]\label{lem:fc-fibre-model}
\notready
\uses{conv:fc-setting, lem:pg-blowup-basic, lem:pg-blowup-chart, lem:fc-affine-fibre-model, lem:fc-fibre-compact, lem:fc-segre-native}
Let $t\in T'^\an$, $x_i$ with $|x_i|_t=1$, and $\theta_t\colon k[x,y]\to\mathcal H(t)[y]$ the substitution
$x_j\mapsto(\bar x_j/\bar x_i)(t)$. Let $Z_t\subseteq\PP^{M'}_{\mathcal H(t)}$ be the closed subscheme defined by
$\theta_t(I')$. There is a homeomorphism $\iota_t\colon\mu^{-1}(t)\to Z_t^\an$ with $|F|_y=|\theta_t(F)|_{\iota_t(y)}$ for every
bihomogeneous $F$. In particular $\mu^{-1}(t)$ is compact. Moreover $Z_t=B_\xi\otimes_{\kappa(\xi)}\mathcal H(t)$, where $\xi$
is the support of $t$ and $B_\xi$ is the scheme-theoretic fibre of $B\to T'$ over $\xi$.
\end{lemma}

\begin{proof}
On each Rees chart $A_i[\mathfrak a_i/g_l]$ (Lemma~\ref{lem:pg-blowup-chart}), the points over $t$ are the
$\mathcal H(t)$-points of $\mathcal H(t)\otimes_{A_i}A_i[\mathfrak a_i/g_l]$, which is the $y_l\ne0$ chart of $Z_t$ (completed
residue field fibre homeomorphism). These identifications agree on overlaps because they are restrictions of one
seminorm. Compactness: $Z_t$ is projective. The last claim is the base change of the fibre.
\end{proof}

\begin{lemma}[Arbitrary psh potentials on Rees charts]\label{lem:fc-chart-psh}
\leanok
\lean{EnvelopeProperty.IsConvexPLOn.fc_comp_completedFibre, EnvelopeProperty.IsDecrConvexPLLimitOn.fc_comp_completedFibre, EnvelopeProperty.IsLocallyPSHOn.fc_comp_completedFibre, EnvelopeProperty.ProjSubvariety.BlowupCenter.fcChartPotential, EnvelopeProperty.ProjSubvariety.BlowupCenter.fcChartPotential_eq_baseTrivialization, EnvelopeProperty.ProjSubvariety.BlowupCenter.fcChartPotential_eq_on_affineFibre, EnvelopeProperty.ProjSubvariety.BlowupCenter.fc_isDecrConvexPLLimitOn_chartPotential, EnvelopeProperty.ProjSubvariety.BlowupCenter.fc_isLocallyPSHOn_chartPotential, EnvelopeProperty.ProjSubvariety.BlowupCenter.fc_isLocallyPSHOn_completedChartFibre, EnvelopeProperty.ProjSubvariety.BlowupCenter.fc_isLocallyPSHOn_nativeChartFibre, EnvelopeProperty.fc_isLocallyPSHOn_affineNativeFibre, EnvelopeProperty.IsLocallyPSHOn.fc_comp_nativeCompletedFibre}
\uses{conv:fc-setting, cor:pg-countable-regularization, lem:pg-blowup-chart, lem:fc-affine-fibre-model}
For every nonzero $g\in\mathfrak a_i$, write $\kappa\colon(\Spec A_i[\mathfrak a_i/g])^\an\to B^\an$ for the Rees
chart map and $\mu_x$ for its first normalization scale. If $\Phi\in\PSH(B,(a,0))$, then
$\Phi\circ\kappa+a\log\mu_x$ is a decreasing limit of convex PL functions, hence locally psh.
Its restriction to the completed-field tensor chart of any fibre is locally psh over that completed field.
On the actual completed scheme fibre of this Rees chart, $\Phi$ itself is intrinsically locally psh.
\end{lemma}

\begin{proof}
Choose decreasing finite Fubini--Study approximants $\Phi_j$ of classes $(a+2^{-j},2^{-j})$.
Reading each approximant on the Rees chart gives a finite maximum $G_j$ of positive rational multiples of
logarithms of regular functions plus rational constants, satisfying
$G_j=\Phi_j\circ\kappa+(a+2^{-j})\log\mu_x+2^{-j}\log\mu_y$.
Both logarithmic scales are nonnegative, so $G_j$ decreases and converges to
$\Phi\circ\kappa+a\log\mu_x$. Taking the maximum with the standard decreasing affine Fubini--Study truncations
provides convex PL functions on the entire chart with the same limit. On the completed-field fibre, replace each
chart coefficient $c$ by $1\otimes c$. Lemma~\ref{lem:fc-affine-fibre-model} identifies all evaluations,
so the same local presentations and decreasing sequences prove local plurisubharmonicity there.
Finite-coefficient affine presentations become presentations in genuine regular units on principal opens; the
actual tensor scheme isomorphism then pulls them back to intrinsic local psh presentations on the same completed
scheme fibre. The normalization term is constant on a fixed fibre because
$\mu_x=\max_j|\bar x_j/\bar x_i|$ depends only on the base point. Subtracting this finite real constant
preserves local plurisubharmonicity and gives the restriction of $\Phi$ itself.
\end{proof}

\begin{lemma}[Restriction to a fibre]\label{lem:fc-fibre-psh}
\notready
\uses{lem:fc-fibre-model, lem:fc-chart-psh, lem:pg-descent-weakly-psh, lem:local-basic}
Let $t\in T'^\an$, and let $(m,s)$, $s\in S_{ma}$, be admissible with $|s|_t\ne0$. Then $m^{-1}\log|s|\circ\mu$ is constant on
$\mu^{-1}(t)$, and $u_t\coloneqq(\Phi-m^{-1}\log|s|\circ\mu)\circ\iota_t^{-1}$ is locally psh on $Z_t^\an$.
\end{lemma}

\begin{proof}
The first claim holds because $|s|_y=|s|_{\mu(y)}$. For the second, the proof of
Lemma~\ref{lem:pg-descent-weakly-psh} applies to $\Phi$ itself, not only to pullbacks: on each Rees chart the dyadic
approximants of $\Phi$ (Corollary~\ref{cor:pg-countable-regularization}) give decreasing convex PL functions in units
($g_m/g_l$, $\bar x_j/\bar x_i$ and $s/x_i^{ma}$). This is \cite[Lemma~5.4(2)]{FHX} on $B$. Pull back along
$Z_t\to B_{\mathcal H(t)}\to B$ (Lemma~\ref{lem:local-basic}(2), Lemma~\ref{lem:fc-base-change}(1)).
\end{proof}

\begin{lemma}[Functions polynomial on lines]\label{lem:fc-lines-poly}
\leanok
\lean{EnvelopeProperty.FcPolynomialOnLines.exists_polynomial_on_open}
Let $F$ be an infinite field, $\Omega\subseteq\A^n_F$ a nonempty open subset, $N\in\N$, and $c\colon\Omega(F)\to F$ a function such
that for every line $L$ meeting $\Omega$ there is $P_L\in F[t]$ of degree $\le N$ with $c=P_L$ on $L\cap\Omega(F)$ (in a
parametrization of $L$). Then $c$ is the restriction of a polynomial.
\end{lemma}

\begin{proof}
Induction on $n$; for $n=1$, $L=\A^1$. For $n\ge2$, choose distinct $a_0,\dots,a_N$ such that the hyperplanes
$H_j=\{w_1=a_j\}$ meet $\Omega$. By induction, $c|_{H_j\cap\Omega}=Q_j(w_2,\dots,w_n)$. Put $Q=\sum_j\ell_j(w_1)Q_j$ with the
Lagrange basis $\ell_j$. If $(a_j,w')\in\Omega$ for all $j$, the line $\{(t,w')\}$ meets $\Omega$, and $P_L$ and $Q(\cdot,w')$
agree at $a_0,\dots,a_N$, so they are equal; hence $c=Q$ on a dense open subset $\Omega'\subseteq\Omega$. For $w\in\Omega(F)$,
a line through $w$ in a direction for which $L\cap\Omega'$ is infinite gives $P_L=Q|_L$, so $c(w)=Q(w)$.
\end{proof}

\begin{lemma}[Affine substitutions and coefficient degrees]\label{lem:fc-factor-degree}
\leanok
\lean{EnvelopeProperty.FcPolynomialOnLines.totalDegree_aeval_affine_le, EnvelopeProperty.FcPolynomialOnLines.natDegree_coeff_factor_bivariateLine_le}
Let $F$ be a field, and regard $P\in F[Z,w_1,\ldots,w_n]$ as a polynomial of total degree $D$.
For $a,b\in F^n$, put $P_L(Z,t)=P(Z,a+tb)$, and suppose $P_L\ne0$.
If $Q\in F[t][Z]$ divides $P_L$, then every coefficient of $Q$ as a polynomial in $Z$ has degree at most $D$ in $t$.
\end{lemma}

\begin{proof}
Substituting affine linear polynomials does not increase total degree. Indeed, a monomial
$c\prod_iw_i^{e_i}$ becomes a polynomial of total degree at most $\sum_i e_i$, and taking the sum over
the support of $P$ gives $\deg P_L\le D$. Write $P_L=QR$ in the domain $F[Z,t]$.
Since $P_L\ne0$, both factors are nonzero, and total degree is additive under multiplication, so
$\deg Q\le\deg P_L\le D$. Each monomial $t^jZ^k$ of $Q$ has $j+k\le\deg Q$;
in particular $j\le D$, which is the claimed coefficient bound.
\end{proof}

\begin{lemma}[Coefficients of factors along lines]\label{lem:fc-factor-coefficients}
\leanok
\lean{EnvelopeProperty.FcPolynomialOnLines.exists_polynomial_of_lineFactor_coeff}
\uses{lem:fc-lines-poly, lem:fc-factor-degree}
Let $F$ be infinite, $\Omega\subseteq\A^n_F$ a nonempty open, $P\in F[w][Z]$ monic, and $k\in\N$.
Suppose $c\colon\Omega(F)\to F$ has the following property: for each affine line $L$ there is a polynomial
$Q_L\in F[t][Z]$ dividing $P|_L$ such that $c|_{L\cap\Omega}$ is the coefficient of $Z^k$ in $Q_L$.
Then $c$ is the restriction of a polynomial in $F[w]$.
\end{lemma}

\begin{proof}
Let $D$ be the total degree of $P$ as a polynomial in $(Z,w)$. Its restriction to every line is nonzero
because it is monic in $Z$. Lemma~\ref{lem:fc-factor-degree} bounds the degree of the coefficient of $Z^k$
in every $Q_L$ by the same number $D$. Lemma~\ref{lem:fc-lines-poly} therefore applies to $c$.
\end{proof}

\begin{lemma}[Globalizing complementary fibre factors]\label{lem:fc-factor-families}
\leanok
\lean{EnvelopeProperty.FcCurveChains.exists_polynomial_factor_family, EnvelopeProperty.FcCurveChains.not_irreducible_of_line_factor_families}
\uses{lem:fc-factor-coefficients}
Let $F$ be infinite, $\Omega\subseteq\A^n_F$ a nonempty open and $P\in F[w][Z]$ monic of $Z$-degree $d$.
Suppose polynomial-valued functions $q_w,q'_w\in F[Z]$ on $\Omega(F)$ satisfy $P_w=q_wq'_w$.
For every $0\le k\le d$ suppose the coefficient functions of $q_w$ and $q'_w$, on every affine line,
are the $Z^k$-coefficients of genuine divisors of $P|_L$ in $F[t][Z]$.
If $q_w$ and $q'_w$ are both nonconstant for some $w\in\Omega(F)$, then $P$ is not irreducible.
\end{lemma}

\begin{proof}
Since $P_w$ is monic and nonzero, divisibility gives $\deg q_w,\deg q'_w\le d$.
Lemma~\ref{lem:fc-factor-coefficients} extends the first $d+1$ coefficient functions of each family to
polynomials in $w$. Assemble them into $R,R'\in F[w][Z]$. The degree bounds show that $R_w=q_w$ and
$R'_w=q'_w$ on $\Omega(F)$, including all coefficients above degree $d$, which vanish.
Thus $P_w=(RR')_w$ there. A polynomial vanishing on a nonempty principal open is zero:
if it vanishes on $D(g)$, then its product with $g$ vanishes at every rational point, hence is the zero
polynomial over the infinite field; the polynomial ring is a domain and $g\ne0$.
Apply this to every coefficient of $P-RR'$ to obtain $P=RR'$.
A unit of $F[w][Z]$ specializes to a unit of $F[Z]$, hence to a constant. The given nonconstant
specializations therefore show that neither $R$ nor $R'$ is a unit, contradicting irreducibility.
\end{proof}

\begin{lemma}[Actual line components through affine points]\label{lem:fc-line-components}
\leanok
\lean{EnvelopeProperty.FcCurveChains.exists_curveStep_over_linePoint, EnvelopeProperty.FcCurveChains.exists_closed_curveStep_over_linePoint}
Let $g\in F[w_1,\ldots,w_n]$ be nonzero, $B$ a domain, and
$F[w]_g\hookrightarrow B$ a finite injective algebra homomorphism. Let $L$ be a nonconstant affine line
such that $g|_L\ne0$. For every scheme point $x\in\Spec B$ over $L\cap D(g)$ and every scheme point
$w\in L\cap D(g)$, there is an irreducible closed curve in $\Spec B$ containing $x$ and a point over $w$.
If $w$ is closed, the point over $w$ can be chosen closed in $\Spec B$.
\end{lemma}

\begin{proof}
The transverse equations of $L$ generate an ideal $J$ with at most $n-1$ generators. Choose a minimal
prime $P$ of $JB$ contained in the prime of $x$. Krull's height theorem, together with $\dim B=n$,
gives $\dim(B/P)>0$. The induced map from the coordinate ring $F[L]_{g|_L}$ to $B/P$ is finite.
Its source is a one-dimensional domain, so a nonzero kernel would force $B/P$ to have dimension zero.
It is therefore injective. The integral extension then gives $\dim(B/P)=1$ and a surjective map
$\Spec(B/P)\to L\cap D(g)$. Choose a point of this curve over $w$; its image in $\Spec B$ and $x$
lie on the same irreducible closed curve $\Spec(B/P)$.
If $w$ is closed, it is still closed after its closed immersion into $D(g)$, and a prime over a maximal
ideal in the finite extension $F[w]_g\hookrightarrow B$ is maximal. Thus the chosen point is closed.
\end{proof}

\begin{lemma}[Actual rational fibres and curve-class factors]\label{lem:fc-rational-fibre-factors}
\leanok
\lean{EnvelopeProperty.FcCurveFibres.rationalFibre_finite, EnvelopeProperty.FcCurveFibres.rationalPoint_isClosed, EnvelopeProperty.FcCurveFibres.rationalPoint_mem_curveClass_iff_on_curve, EnvelopeProperty.FcCurveFibres.classFactor_monic, EnvelopeProperty.FcCurveFibres.classFactor_mul_compl, EnvelopeProperty.FcCurveFibres.natDegree_classFactor, EnvelopeProperty.FcCurveFibres.natDegree_classFactor_pos_of_mem}
Let $\varphi\colon R\to B$ be a finite $F$-algebra homomorphism and $w\colon R\to F$ an $F$-rational point.
Then the actual rational fibre $s_w=\{\chi\colon B\to F:\chi\varphi=w\}$ is finite, and its underlying scheme
points are closed. For a curve class $S$, put
$q_{S,w}=\prod_{\chi\in s_w,\,\ker\chi\in S}(Z-\chi(b))$.
These factors are monic, $q_{S,w}q_{S^c,w}=\prod_{\chi\in s_w}(Z-\chi(b))$, and
$\deg q_{S,w}=\#\{\chi\in s_w:\ker\chi\in S\}$. Thus the factor is nonconstant if its selected subset is nonempty.
Membership in $S$ is constant on the rational points of any actual irreducible closed affine curve.
\end{lemma}

\begin{proof}
Give $B$ and $F$ their $R$-algebra structures through $\varphi$ and $w$. A homomorphism in the actual rational fibre
is an $R$-algebra homomorphism $B\to F$, and these form a finite set because $B$ is a finite $R$-module and $F$ is a field.
The map from the actual fibre to this finite set is injective on the entire algebra homomorphism.
Each $F$-algebra homomorphism to $F$
is surjective because it maps the scalar $a$ to $a$, so its kernel is maximal and its scheme point is closed.
The product identity partitions a finite product into the selected subset and its complement. All linear factors
are monic of degree one, giving monicity and the exact degree formula, including repeated values.
Finally, any two points on the same actual irreducible closed curve are joined by a curve step; appending that
step to a chain, in either direction, proves constancy of curve-class membership.
\end{proof}

\begin{lemma}[Rational points on actual line components]\label{lem:fc-rational-line-components}
\leanok
\lean{EnvelopeProperty.FcCurveFibres.exists_rational_curveStep_over_linePoint}
\uses{lem:fc-line-components, lem:fc-rational-fibre-factors}
Let $F$ be algebraically closed and $F[w]_g\hookrightarrow B$ a finite injective homomorphism to a domain.
Let $L$ be a nonconstant affine line with $g|_L\ne0$, let $x\colon B\to F$ lie over $L\cap D(g)$, and let
$w\colon F[L]_{g|_L}\to F$ be a rational point of the line. Then there is an actual rational point
$y\colon B\to F$ joined to $x$ by one irreducible closed curve, such that the entire restriction of $y$ to
$F[w]_g$ equals the restriction of $w$ along the line homomorphism.
\end{lemma}

\begin{proof}
Lemma~\ref{lem:fc-line-components} supplies a closed point of $\Spec B$ over $w$, joined to $x$ by an actual curve.
The finite map makes $B$ an algebra of finite type over $F$. By Nullstellensatz the residue field of this closed
point is $F$, so composing the quotient map with the inverse scalar isomorphism gives an actual algebra
homomorphism $y\colon B\to F$. Its base restriction has the same kernel as that of $w$.
Two $F$-algebra homomorphisms to $F$ with the same kernel agree: for every $r$, subtract the scalar which is the
value of the first homomorphism, and apply the second. This proves equality of the entire base homomorphisms.
\end{proof}

\begin{lemma}[Every actual rational curve class meets every fibre]\label{lem:fc-classes-meet-fibres}
\leanok
\lean{EnvelopeProperty.FcCurveFibres.exists_rational_curveChain_in_fibre, EnvelopeProperty.FcCurveFibres.natDegree_fibreCurveClassFactor_pos, EnvelopeProperty.FcCurveFibres.natDegree_fibreCurveClassFactor_compl_pos}
\uses{lem:fc-rational-line-components, lem:fc-rational-fibre-factors}
Let $F$ be algebraically closed and $F[w]_g\hookrightarrow B$ a finite injective homomorphism to a domain.
Every curve class containing an actual rational point of $\Spec B$ meets every rational fibre over $D(g)$.
Its root-product factor is therefore nonconstant in every such fibre. If another rational point lies outside
the class, the complementary factor is also nonconstant in every rational fibre.
\end{lemma}

\begin{proof}
For a rational point $x$ and a prescribed base point $w$, compare their entire base homomorphisms.
If these agree, use $x$ itself and the empty chain. Otherwise their coordinate tuples $a,z$ differ, so
$b=z-a$ has a nonzero coordinate. The genuine affine line $a+tb$ passes through $a$ at zero and $z$ at one.
The value $g(a)$ is nonzero because $g$ is a unit on the principal open, so $g$ restricts to a nonzero polynomial
on this line. The line ideal vanishes at $a$, and $g(z)\ne0$ supplies the actual rational line point at one.
Lemma~\ref{lem:fc-rational-line-components} constructs a rational point over the whole prescribed base point,
joined to $x$ by a genuine curve. Thus it is in the same class. The exact root-product degree formula gives
positive degree. Apply the same argument starting from a point outside the class; symmetry and concatenation
of chains show that its class is disjoint from the first class, proving positivity for the complement.
\end{proof}

\begin{lemma}[Identifying actual monogenic fibre factors]\label{lem:fc-monogenic-fibre-factors}
\leanok
\lean{EnvelopeProperty.FcMonogenicFibres.rationalFibreEquivRoots, EnvelopeProperty.FcMonogenicFibres.root_injective_on_rationalFibre, EnvelopeProperty.FcMonogenicFibres.image_root_rationalFibreFinset, EnvelopeProperty.FcMonogenicFibres.prod_rationalFibre_eq_map, EnvelopeProperty.FcMonogenicFibres.fibreClassFactor_mul_compl_eq_map, EnvelopeProperty.FcMonogenicFibres.fibreClassFactor_dvd_map}
\uses{lem:fc-rational-fibre-factors}
Let $P\in R[Z]$ be monic, where $R$ is an algebra over an algebraically closed field $F$, and let
$w\colon R\to F$ be an actual rational base point such that $P_w$ is separable.
The actual rational fibre of $R[Z]/(P)$ is in bijection with the roots of $P_w$, via evaluation of the adjoined root.
Consequently, the two complementary actual root products attached to any subset of its scheme points multiply
to $P_w$, and each product is a genuine divisor of $P_w$.
\end{lemma}

\begin{proof}
Any rational homomorphism of the quotient evaluates its adjoined root to a root of $P_w$.
Conversely a root $z$ of $P_w$ induces the whole rational homomorphism by $r\mapsto w(r)$ and $Z\mapsto z$;
the polynomial quotient relation holds precisely because $P_w(z)=0$. Two such homomorphisms agree when their
whole base restrictions and root values agree. Monicity makes the quotient a finite $R$-module, so its actual
rational fibre is finite. The image of root evaluation on that finite fibre is exactly the root set of $P_w$.
Since $F$ is algebraically closed and $P_w$ is separable and monic, every root occurs exactly once in its complete
linear factorization. Partition this actual fibre into the chosen subset and its complement.
\end{proof}

\begin{lemma}[Actual curves of irreducible line factors]\label{lem:fc-line-factor-curves}
\leanok
\lean{EnvelopeProperty.FcLineFactorCurves.adjoinRoot_map_surjective, EnvelopeProperty.FcLineFactorCurves.ringKrullDim_adjoinRoot_line_factor, EnvelopeProperty.FcLineFactorCurves.curveStep_of_irreducible_line_factor, EnvelopeProperty.FcLineFactorCurves.exists_rationalPoint_on_line_factor, EnvelopeProperty.FcLineFactorCurves.curveClass_iff_lineFactorPoint}
\uses{lem:fc-monogenic-fibre-factors, lem:fc-line-components}
Let $P\in F[w][Z]$, let $L\colon w=a+tb$ be a genuine nonconstant affine line, and let $Q\in F[t][Z]$
be a monic irreducible factor of $P|_L$. Any two actual rational points of $F[w][Z]/(P)$ lying on the
specializations of $Q$ belong to one actual irreducible closed curve. If $F$ is algebraically closed, $Q$
has an actual rational point over each line parameter, and membership in any curve class is determined by
one constructed reference point over parameter zero.
\end{lemma}

\begin{proof}
The actual substitution $F[w]\to F[t]$ is surjective because $b$ has a nonzero coordinate. It induces a
surjective homomorphism $F[w][Z]/(P)\to F[t][Z]/(Q)$, by lifting polynomial representatives coefficientwise.
The target is a domain because $Q$ is irreducible over the unique factorization domain $F[t]$.
Monicity makes this quotient finite over $F[t]$, and its positive degree makes the scalar map injective.
Normality of $F[t]$ and integrality therefore give dimension one. The kernel is consequently a prime ideal
defining an actual irreducible closed curve in the original spectrum. Each specialized root of $Q$ constructs
an entire homomorphism of this quotient; agreement on all base coefficients and the adjoined root proves
that its composite is the prescribed original rational point. Thus both prescribed points lie on this curve.
Over an algebraically closed field, specialization preserves the positive degree of the monic polynomial
$Q$, so every specialization has a root. Choose one at parameter zero. The constructed curve step shows
that every other specialized root has the same curve-class membership as that reference point.
\end{proof}

\begin{lemma}[Actual polynomial line factors of monogenic curve classes]\label{lem:fc-monogenic-line-factors}
\leanok
\lean{EnvelopeProperty.FcMonogenicLineFactors.exists_monic_irreducible_factorization, EnvelopeProperty.FcMonogenicLineFactors.exists_curveClass_lineFactor, EnvelopeProperty.FcMonogenicLineFactors.exists_curveClass_complementary_lineFactors}
\uses{lem:fc-monogenic-fibre-factors, lem:fc-line-factor-curves}
Let $F$ be algebraically closed, $P\in F[w][Z]$ monic, $S$ an actual curve class in $\Spec F[w][Z]/(P)$,
and $L\colon w=a+tb$ a nonconstant actual affine line. There are complementary monic polynomial factors
$D,E\in F[t][Z]$ with $DE=P|_L$, such that at every parameter where $P(a+tb,Z)$ is separable, $D_t$ is the
actual rational root product of $S$ and $E_t$ is the actual rational root product of its complement.
\end{lemma}

\begin{proof}
Degree induction factors any monic polynomial over a domain into finitely many monic irreducibles:
if the polynomial is neither one nor irreducible, the monic irreducibility criterion gives two nonconstant
monic factors of smaller degree. Apply this to $P|_L=\prod_iQ_i$.
For each $Q_i$, Lemma~\ref{lem:fc-line-factor-curves} constructs an actual reference point over parameter zero.
Let $D$ be the product of those entire $Q_i$ whose reference point lies in $S$. Then $D$ is monic and divides
$P|_L$. At any parameter, a root of $D_t$ belongs to one of the selected factors, and hence lies in $S$ by the
actual curve joining it to the reference point. Conversely a point of $S$ is a root of $P_t$, so is a root of
some $Q_{i,t}$; the same curve construction shows that this factor was selected. Thus $D_t$ and the actual
class factor have exactly the same roots. Both divide the separable monic polynomial $P_t$, so both are
separable and monic; over the algebraically closed field this root equality proves polynomial equality.
Write $P|_L=DE$. The complementary factor $E$ is monic, and the actual whole-fibre factorization of
Lemma~\ref{lem:fc-monogenic-fibre-factors}, together with cancellation of the nonzero monic $D_t$, identifies
$E_t$ with the actual complementary class factor.
\end{proof}

\begin{lemma}[Actual curve closures in the original finite algebra]\label{lem:fc-birational-curve-closures}
\leanok
\lean{EnvelopeProperty.FcCurveClosures.kernel_quotient_dim_one_of_line_localization, EnvelopeProperty.FcCurveClosures.commonOpenCurveMap, EnvelopeProperty.FcCurveClosures.commonOpen_rationalHom_ext, EnvelopeProperty.FcCurveClosures.curveStep_of_commonOpen_irreducible_line_factor}
\uses{lem:fc-line-factor-curves}
Let $C=F[w][Z]/(P)$, let $\delta\colon C\to A$ have finite composite $F[w]\to A$, and suppose that
$C_h\simeq A_{\delta(h)}$ is a genuine isomorphism compatible with the whole homomorphism $\delta$.
Let $Q\in F[t][Z]$ be a monic irreducible factor of the restriction of $P$ to a nonconstant actual affine line.
Two prescribed actual rational points of $\Spec A$ in this common open whose base and generator values lie on
the specializations of $Q$ are joined by an actual irreducible closed curve in $\Spec A$.
\end{lemma}

\begin{proof}
Put $D=F[t][Z]/(Q)$ and let $\kappa\colon C\to D$ be the actual polynomial quotient map. Construct
$\alpha\colon A\to A_{\delta(h)}\simeq C_h\to D_{\kappa(h)}$.
The prescribed rational point shows $\kappa(h)\ne0$. Thus the target is a domain, and $J=\ker\alpha$ is prime.
The whole base diagram sends the transverse line equations to zero in $A/J$, so the surjective line substitution
induces $F[t]\to A/J$. The original finiteness over $F[w]$ descends by restriction of scalars to finiteness
over $F[t]$. Composing with $A/J\hookrightarrow D_{\kappa(h)}$ proves faithfulness because the monic irreducible
polynomial $Q$ has positive degree and makes $F[t]\hookrightarrow D$ injective. Normality and dimension one of
$F[t]$ therefore give $\dim(A/J)=1$.
The prescribed generator values construct rational homomorphisms $D\to F$. Their entire restrictions to $C$
equal those of the prescribed original points; the nonvanishing of $h$ extends them to $D_{\kappa(h)}$.
On the common open the restriction to $C$ determines the entire rational homomorphism of $A$: extend both
homomorphisms to $A_{\delta(h)}$, compare them through $C_h$, and use surjectivity of the actual isomorphism.
Thus their composites through $\alpha$ are exactly the prescribed original points, so both contain the actual
prime $J$ and lie on the genuine closed curve $\Spec(A/J)$ in the original spectrum.
\end{proof}

\begin{lemma}[Constructing the finite birational common open]\label{lem:fc-finite-birational-open}
\leanok
\lean{EnvelopeProperty.FcFiniteBirationalOpen.exists_common_denominator, EnvelopeProperty.FcFiniteBirationalOpen.exists_commonPrincipalOpen_equiv, EnvelopeProperty.FcFiniteBirationalOpen.exists_commonPrincipalOpen_with_base_denominator}
\uses{lem:nf-charpoly}
Let $C$ be a domain and $A$ a finite $C$-algebra with an injective $C$-algebra homomorphism into
$\operatorname{Frac} C$. There is a constructed nonzero $h\in C$ such that the canonical map
$C_h\to A_h$ is an actual isomorphism, compatible with the entire original map $C\to A$.
If $C$ is finite and faithful over a normal domain $R$, there is also a constructed $g\in R\setminus\{0\}$
such that $h$ divides the image of $g$ in $C$. Thus this common open contains every rational fibre over $D(g)$.
\end{lemma}

\begin{proof}
Choose finitely many $C$-module generators of $A$. Clearing the denominators of their images in
$\operatorname{Frac} C$ gives one nonzero $h$. The property that $h$ times an element has image in $C$
is preserved by addition and scalar multiplication, so span induction gives it for every element of $A$.
Injectivity into the fraction field identifies these equalities in the original algebra itself:
for every $a\in A$, $ha$ is the image of some $c\in C$. The canonical localized map is injective because
$C\to A$ is injective, and this uniform denominator proves it is surjective; hence it is an actual algebra
isomorphism. Its localization formula retains all original coefficient values.
For the final assertion take $g$ to be the constant coefficient of the characteristic polynomial of $h$
over $R$. Lemma~\ref{lem:nf-charpoly} proves both $g\ne0$ and $h\mid g$ in $C$. A rational point over $D(g)$
therefore has nonzero value on $h$, proving the asserted whole-fibre inclusion.
\end{proof}

\begin{lemma}[Constructing a global separating normalization]\label{lem:fc-separating-normalization}
\leanok
\lean{EnvelopeProperty.FcSeparatingNormalization.exists_finite_separating_normalization_with_fractionMap, EnvelopeProperty.FcSeparatingNormalization.exists_finite_injective_genericallyEtale_normalization}
\uses{lem:projection-existence}
Every finite-type domain $A$ over a perfect field $F$ admits a constructed global finite injective map
$F[w_1,\ldots,w_n]\hookrightarrow A$ that is generically \'etale. Its actual compatible fraction-field map is
also constructed, and the fraction-field extension is formally \'etale. No normalization, dimension, or
generic-separability hypothesis is supplied.
\end{lemma}

\begin{proof}
Start with Noether normalization by coordinates $\tau_1,\ldots,\tau_n$. The differential perturbation
construction used in Lemma~\ref{lem:projection-existence}, applied to the whole nonzero ideal $A$ and to $n$
zero elements, supplies $a_1,\ldots,a_n\in A$ whose differentials span or are independent in
$\Omega_{\operatorname{Frac}A/F}$. Over a perfect field that differential space has dimension at most $n$,
so they span. Apply the integral-coordinate twist to all coordinates, choosing sufficiently large exponents
$m_i$ divisible by the exponential characteristic. Then $y_i=\tau_i^{m_i}-a_i$ still make $A$ integral over
$F[y]$. Dimension $n$ proves the resulting polynomial map injective, and finite type plus integrality makes it finite.
In characteristic zero generic separability is automatic. In characteristic $p$, $p\mid m_i$ gives
$dy_i=-da_i$, so these new coordinate differentials still span. The actual fraction-field homomorphism is
constructed by localization of the injective polynomial map; the differential criterion proves it formally
\'etale. Openness of the \'etale locus for the finite map constructs a nonzero principal open on which the
actual localized map is \'etale.
\end{proof}

\begin{lemma}[Constructing an integral primitive element]\label{lem:fc-integral-primitive-element}
\leanok
\lean{EnvelopeProperty.FcPrimitiveIntegralGenerator.exists_primitive_element_mem_ring, EnvelopeProperty.FcPrimitiveIntegralGenerator.exists_global_normalization_with_integral_primitive_element}
\uses{lem:fc-separating-normalization}
For every finite-type domain $A$ over a perfect field $F$, one can construct a global finite injective
generically \'etale normalization $R=F[w]\hookrightarrow A$ and an element $b\in A$ such that
$\operatorname{Frac}A=\operatorname{Frac}R(b)$. Neither normalization nor primitive element is an input.
The original algebra is not required to be normal.
\end{lemma}

\begin{proof}
Lemma~\ref{lem:fc-separating-normalization} constructs the actual compatible formally \'etale fraction-field
map $K=\operatorname{Frac}R\to L=\operatorname{Frac}A$. Finite type and formal unramifiedness over the field
make this a finite separable extension, so choose a field primitive element $\alpha\in L$.
Write $\alpha=a/s$ with $a,s\in A$, $s\ne0$. Algebraicity of $s$ over $R$ gives $b\in A$ and
$d\in R\setminus\{0\}$ with $da=sb$. Hence $d\alpha=b$ in the actual fraction field. Since $d$ is an invertible
scalar from $K$, $K(b)=K(\alpha)=L$. Thus the generator lies in the original finite algebra and is integral over $R$.
\end{proof}

\begin{lemma}[Constructing the separable specialization locus]\label{lem:fc-separable-good-locus}
\leanok
\lean{EnvelopeProperty.FcSeparableGoodLocus.exists_nonzero_bezout_derivative, EnvelopeProperty.FcSeparableGoodLocus.exists_separable_principal_locus}
Let $R$ be a domain with fraction field $K$, and let $P\in R[Z]$ have separable image in $K[Z]$.
There is a constructed $g\in R\setminus\{0\}$ such that for every homomorphism $w\colon R\to E$ to a field,
$w(g)\ne0$ implies that the entire specialized polynomial $P_w$ is separable.
\end{lemma}

\begin{proof}
Generic separability gives $UP+VP'=1$ in $K[Z]$. Clear the coefficients of $U$ and $V$ to obtain
$U_0,V_0\in R[Z]$ and nonzero $r,s\in R$ with $U_0=rU$ and $V_0=sV$ in the fraction field.
Then $(sU_0)P+(rV_0)P'=rs$ holds in $R[Z]$, by injectivity into $K[Z]$.
Take $g=rs$. If $w(g)\ne0$, apply the whole polynomial homomorphism induced by $w$ and divide this identity
by the nonzero scalar $w(g)$. This is a B\'ezout identity for $P_w$ and its derivative, proving separability.
\end{proof}

\begin{lemma}[The actual primitive monogenic subalgebra]\label{lem:fc-primitive-monogenic-subalgebra}
\leanok
\lean{EnvelopeProperty.FcPrimitiveMonogenicSubalgebra.exists_fraction_embedding_of_primitive_adjoin, EnvelopeProperty.FcPrimitiveMonogenicSubalgebra.exists_commonOpen_of_primitive_adjoin, EnvelopeProperty.FcPrimitiveMonogenicSubalgebra.monic_irreducible_minpoly_presentation, EnvelopeProperty.FcPrimitiveMonogenicSubalgebra.commonOpenAlgEquiv_of_ringEquiv}
\uses{lem:fc-integral-primitive-element, lem:fc-finite-birational-open}
Let $R\hookrightarrow A$ be a finite inclusion of domains with fraction fields $K\subseteq L$, and let
$b\in A$ satisfy $K(b)=L$. For the actual subalgebra $C=R[b]\subseteq A$, the original algebra embeds into
$\operatorname{Frac} C$ as a $C$-algebra. There are constructed nonzero $h\in C$ and $g\in R$ with $h\mid g$
in $C$ and an actual isomorphism $C_h\simeq A_h$ preserving the whole original map $C\to A$.
If $R$ is normal, $C$ has the actual presentation $R[Z]/(P)$ for the monic irreducible polynomial
$P=\minpoly_R(b)$, and the adjoined root maps to the prescribed original element $b$.
\end{lemma}

\begin{proof}
Construct the canonical injective map $\sigma\colon\operatorname{Frac}C\to L$ and the compatible scalar map
$K\to\operatorname{Frac}C$ by localization of the whole original ring maps. The field range of $\sigma$
contains $K$ and the actual element $b$, so $K(b)=L$ proves that $\sigma$ is surjective. Its inverse, composed
with $A\hookrightarrow L$, is the required entire $C$-algebra embedding of $A$ into $\operatorname{Frac}C$.
The singleton adjoin is finite over $R$ because $b$ is integral, and $A$ is finite over $C$ by restriction of
scalars. Lemma~\ref{lem:fc-finite-birational-open} constructs $h$ and the canonical common-open isomorphism.
Algebraicity of $h$ over $R$ supplies a nonzero $g\in R$ divisible by $h$ in $C$.
For normal $R$, the actual minimal-polynomial quotient isomorphism identifies $R[Z]/(P)$ with the singleton
adjoin, retaining the actual root value. Finally, compatibility with every original coefficient upgrades the
constructed common-open ring isomorphism to an algebra isomorphism over any existing coefficient field.
\end{proof}

\begin{lemma}[The constructed monogenic birational good locus]\label{lem:fc-monogenic-birational-presentation}
\leanok
\lean{EnvelopeProperty.FcMonogenicBirationalPresentation.exists_minpoly_presentation_commonOpen, EnvelopeProperty.FcMonogenicBirationalPresentation.exists_minpoly_presentation_goodLocus, EnvelopeProperty.FcMonogenicBirationalPresentation.exists_global_monic_presentation_goodLocus}
\uses{lem:fc-primitive-monogenic-subalgebra, lem:fc-separable-good-locus}
For every finite-type domain $A$ over a perfect field $F$, one can construct a finite injective normalization
$\varphi\colon R=F[w]\to A$, a monic irreducible $P\in R[Z]$, an injective actual map
$\delta\colon R[Z]/(P)\to A$ whose entire restriction to $R$ is $\varphi$, and nonzero $h\in R[Z]/(P)$,
$g\in R$ with $h\mid g$, together with an actual $F$-algebra isomorphism
$(R[Z]/(P))_h\simeq A_{\delta(h)}$ preserving every original quotient coefficient.
For every actual rational base homomorphism $w\colon R\to F$ with $w(g)\ne0$, $P_w$ is separable.
No normalization, polynomial presentation, primitive element, common open, or good locus is an input.
\end{lemma}

\begin{proof}
Construct the actual primitive element and singleton adjoin by the preceding lemmas. Present the actual
subalgebra by $P=\minpoly_R(b)$, with its prescribed root value. Pull the constructed conductor $h$ back
through this genuine quotient isomorphism, and transport the actual common-open isomorphism by localization
of that entire ring isomorphism. Applying the inverse polynomial quotient map to $h\mid g$ preserves the
original nonzero base denominator and all coefficient values.
Over the fraction field of the normal base, $P$ is the actual minimal polynomial of the prescribed element.
The formally \'etale fraction extension makes this polynomial separable. Lemma~\ref{lem:fc-separable-good-locus}
constructs a second nonzero base element controlling every separable specialization. Multiply the two
base elements: divisibility by the conductor persists, and nonvanishing of the product implies the required
specialized separability. Finally upgrade the whole-coefficient-compatible common-open ring isomorphism to
the actual $F$-algebra isomorphism. Each map and open in the statement is therefore constructed from $A$ itself.
\end{proof}

\begin{lemma}[Identifying actual original-algebra fibre factors]\label{lem:fc-original-fibre-factors}
\leanok
\lean{EnvelopeProperty.FcCommonOpenFibres.commonOpenRationalExtension_comp, EnvelopeProperty.FcCommonOpenFibres.rationalFibreEquivRoots_of_commonOpen, EnvelopeProperty.FcCommonOpenFibres.root_injective_on_rationalFibre_of_commonOpen, EnvelopeProperty.FcCommonOpenFibres.prod_rationalFibre_eq_map_of_commonOpen, EnvelopeProperty.FcCommonOpenFibres.fibreClassFactor_mul_compl_eq_map_of_commonOpen, EnvelopeProperty.FcCommonOpenFibres.fibreClassFactor_dvd_map_of_commonOpen}
\uses{lem:fc-monogenic-birational-presentation, lem:fc-rational-fibre-factors, lem:fc-birational-curve-closures}
In the actual presentation and common open of Lemma~\ref{lem:fc-monogenic-birational-presentation}, the rational
fibre of the ORIGINAL algebra $A$ over every good base homomorphism $w$ is in bijection with the roots of $P_w$,
by evaluating $\delta(Z)$. This element separates that entire actual fibre. If $F$ is algebraically closed and
$P_w$ is separable, $P_w$ is exactly the root product over the actual rational points of $A$ above $w$.
Consequently the two complementary root products for any actual subset of $\Spec A$ multiply to $P_w$,
and each is a genuine polynomial divisor of $P_w$.
\end{lemma}

\begin{proof}
The divisibility $h\mid g$ and $w(g)\ne0$ force every rational homomorphism over $w$ to be nonzero on $h$.
A rational homomorphism of $A$ restricts to the actual polynomial quotient, hence evaluates its root to a root
of $P_w$. Conversely such a root constructs the entire quotient homomorphism, which extends to its $h$-localization.
Compose with the inverse of the genuine common-open isomorphism and the original localization map of $A$.
This constructs an actual homomorphism $A\to F$ and retains its WHOLE prescribed quotient restriction.
Uniqueness follows because homomorphisms on the common open are determined by that entire restriction;
the original points are therefore identified exactly. This proves the actual fibre/root bijection and separation.
The finite original map makes this actual fibre finite. Over an algebraically closed field, the monic separable
polynomial has every root exactly once; the bijection therefore identifies its entire product of linear factors
with the actual original-algebra fibre root product. Partitioning that finite product by an actual subset of
$\Spec A$ and its complement gives the final assertions.
\end{proof}

\begin{lemma}[Actual original-algebra polynomial line factors]\label{lem:fc-original-line-factors}
\leanok
\lean{EnvelopeProperty.FcCommonOpenLineFactors.exists_commonOpen_point_on_line_factor, EnvelopeProperty.FcCommonOpenLineFactors.exists_curveClass_lineFactor_of_commonOpen, EnvelopeProperty.FcCommonOpenLineFactors.exists_curveClass_complementary_lineFactors_of_commonOpen}
\uses{lem:fc-original-fibre-factors, lem:fc-birational-curve-closures, lem:fc-monogenic-line-factors}
In the constructed common-open presentation of the ORIGINAL finite algebra $A$, let $S$ be an actual curve
class in $\Spec A$ and let $L\colon w=a+tb$ be a genuine nonconstant affine line.
There are constructed complementary monic polynomial factors $D,E\in F[t][Z]$, $DE=P|_L$, whose
specializations on every good separable fibre are the actual ORIGINAL-algebra root products of $S$ and its complement.
No line divisor, class grouping, chosen reference point, or identification of original points is an input.
\end{lemma}

\begin{proof}
If $g|_L=0$, there are no good parameters and take $D=1$, $E=P|_L$.
Otherwise choose $t_0$ with $g(a+t_0b)\ne0$ and construct a finite monic irreducible factorization
$P|_L=\prod_iQ_i$. For each factor, a specialized root at $t_0$ constructs an actual rational point of the
monogenic quotient. Since $h\mid g$ it lies on the genuine common open, and extending through its actual
isomorphism constructs an ORIGINAL-algebra point $x_i\colon A\to F$, preserving the entire quotient base
restriction, generator value and conductor nonvanishing. Select the entire $Q_i$ precisely when $\ker x_i\in S$,
and let $D$ be their product. Then $D$ is monic and divides $P|_L$.
For any original rational point over a good parameter whose generator value is a root of $Q_i$,
Lemma~\ref{lem:fc-birational-curve-closures} constructs a genuine irreducible CLOSED curve in $\Spec A$
joining that exact original point to $x_i$. Thus its membership in the original class $S$ is exactly the
selection of $Q_i$. This uses actual curves in $A$, not curves in the monogenic subalgebra treated as curves in $A$.
Every original rational fibre is identified with the specialized roots by
Lemma~\ref{lem:fc-original-fibre-factors}. Consequently $D_t$ and the actual original-class factor have exactly
the same roots. Both are monic divisors of the separable $P_t$, so their roots occur once and polynomial
equality follows. The monic complementary polynomial $E=P|_L/D$ specializes to the actual complementary
factor by the whole original-fibre product identity and cancellation of the nonzero monic $D_t$.
\end{proof}

\begin{lemma}[Unconditional affine integral curve chains]\label{lem:fc-affine-integral-curve-chains}
\leanok
\lean{EnvelopeProperty.FcAffineCurveChains.exists_curveChain_in_global_fibre, EnvelopeProperty.FcAffineCurveChains.curveChain_of_finiteType_domain, EnvelopeProperty.FcAffineCurveChains.curveChain_of_isMaximal, EnvelopeProperty.FcAffineCurveChains.curveChain_of_isClosed}
\uses{lem:fc-factor-families, lem:fc-classes-meet-fibres, lem:fc-monogenic-birational-presentation, lem:fc-original-fibre-factors, lem:fc-original-line-factors}
Let $F$ be algebraically closed and $A$ a finite-type integral $F$-algebra. Any two actual rational
points of $\Spec A$, and hence any two closed points, are joined by a finite chain of genuine
irreducible closed curves in $\Spec A$. No normalization, primitive element, separability,
good locus, line factor, class count or geometric lifting assumption is an input.
\end{lemma}

\begin{proof}
First, for any finite injective polynomial normalization, every actual curve class meets every rational
fibre: apply Lemma~\ref{lem:fc-classes-meet-fibres} through the actual localization isomorphism at $1$.
Construct the finite injective polynomial normalization $\varphi\colon R=F[w]\to A$, the monic
irreducible polynomial $P\in R[Z]$, its actual quotient map into $A$, and the common monogenic good
locus $D(g)$ using Lemma~\ref{lem:fc-monogenic-birational-presentation}. Choose a rational base point
$u\in D(g)$ by polynomial nonvanishing over the infinite field $F$.

Suppose two original rational points $x,y$ belong to distinct curve classes, and let $S$ be the class of $x$.
At every base point define $q$ and $q'$ by the actual finite rational fibre root products for $S$ and
its complement. On $D(g)$, Lemma~\ref{lem:fc-original-fibre-factors} gives $P_w=q_wq'_w$.
On every nonconstant affine line, Lemma~\ref{lem:fc-original-line-factors} constructs genuine
complementary polynomial line factors whose entire specializations on $D(g)$ are $q$ and $q'$.
Taking coefficients therefore gives exactly the line-polynomial conditions of
Lemma~\ref{lem:fc-factor-families}; no presumed polynomial coefficient functions are introduced.
Both classes meet the actual fibre over $u$, so the degrees of both factors there are strictly positive.
The complementary point remains outside $S$, since otherwise its chain to $y$ would join $x$ and $y$.
Lemma~\ref{lem:fc-factor-families} now contradicts irreducibility of the constructed $P$.
Thus the rational points are joined. Nullstellensatz constructs actual rational homomorphisms for
any prescribed closed points, proving the closed-point form as well. Every curve used here is an
actual prime quotient of $A$ of Krull dimension one, as required by the curve-step relation.
\end{proof}

\begin{lemma}[Connected affine component reduction]\label{lem:fc-connected-affine-curve-chains}
\leanok
\lean{EnvelopeProperty.FcAffineCurveChains.curveChain_of_isClosed}
\uses{lem:fc-affine-integral-curve-chains}
Let $F$ be algebraically closed and $A$ any finite-type $F$-algebra with connected spectrum.
Any two closed points of $\Spec A$ are joined by genuine irreducible closed curves in that original spectrum.
No reducedness, irreducibility, component cover, or intersection point is an input.
\end{lemma}

\begin{proof}
A surjective map $A\to B$ transports every actual curve step: if $Q\subseteq B$ is the prime
of the curve, then $A/f^{-1}(Q)\simeq B/Q$ retains Krull dimension one and both endpoint containments.
Transport the entire finite chain step by step. In particular, any two original closed points containing
an actual prime $P\subseteq A$ lift to closed points of the domain $A/P$, and
Lemma~\ref{lem:fc-affine-integral-curve-chains} joins them by actual original-spectrum curves.

Fix one closed point $x$. A minimal prime is good if its closed subset contains a closed point
joined to $x$. The union $U$ of the good minimal-prime closed subsets and the union $V$ of the
remaining ones are closed, since the finite-type algebra is noetherian and has finitely many minimal primes.
They cover the entire spectrum. They are disjoint: if a point lies in a good and a remaining
component, choose an actual maximal ideal containing that point. The preceding common-prime result joins
this closed specialization to the chosen good-component witness. It therefore witnesses goodness of
the remaining component as well, a contradiction. Thus $U$ is clopen. It contains $x$, so connectedness
forces $U$ to be the entire spectrum. Every other closed point belongs to a good component, and the
common-prime theorem joins it to that component's actual witness and hence to $x$.
\end{proof}

\begin{lemma}[Dimensions of actual scheme curve closures]\label{lem:fc-scheme-curve-dimensions}
\leanok
\lean{EnvelopeProperty.FcSchemeCurveDimensions.topologicalKrullDim_le_one_of_open_neighbourhoods, EnvelopeProperty.FcSchemeCurveDimensions.ringKrullDim_basicOpen, EnvelopeProperty.FcSchemeCurveDimensions.ringKrullDim_affine_eq, EnvelopeProperty.FcSchemeCurveDimensions.topologicalKrullDim_affine, EnvelopeProperty.FcSchemeCurveDimensions.topologicalKrullDim_eq_one_of_affine, EnvelopeProperty.FcSchemeCurveDimensions.isReduced_vanishingIdeal_subscheme, EnvelopeProperty.FcSchemeCurveDimensions.isIntegral_vanishingIdeal_subscheme, EnvelopeProperty.FcSchemeCurveDimensions.exists_closed_integral_curve_of_affine_prime}
\uses{lem:fc-connected-affine-curve-chains}
In an integral scheme locally of finite type over a field, every nonempty affine coordinate ring has
the same dimension. If one has dimension one, the entire actual scheme has topological Krull dimension one.
The reduced closed closure of an actual irreducible affine curve in a finite-type scheme is therefore
an actual integral closed curve of dimension one in the original scheme.
\end{lemma}

\begin{proof}
Dimension at most one is local on open neighbourhoods: for a strict three-term chain of irreducible
closed sets, choose an open neighbourhood of a point in the smallest member. Restriction is an order
isomorphism for irreducible closed subsets meeting that open, so all three strict steps remain and
contradict its dimension bound. In an integral finite-type scheme, two nonempty affine opens meet and
have a common principal affine open. Each actual coordinate ring of that principal open is the actual
localization of the corresponding affine coordinate ring. Finite-type integral principal localization
preserves dimension, so the two original affine dimensions are equal. Thus one affine curve chart gives
dimension at most one on every affine neighbourhood, and its subspace dimension gives the reverse inequality.

For an actual affine prime curve, close its image in the original scheme and equip that closed set with
its actual vanishing-ideal subscheme. On every affine chart the defining ideal is radical, so its actual
quotient is reduced, including empty intersections. Irreducibility of the image and its closure therefore
make the constructed whole subscheme integral. The preimage of the closure in the original affine open is
exactly the original closed curve, since the affine inclusion is an open embedding. Consequently the
constructed closed subscheme has an actual affine coordinate quotient of dimension one, and the preceding
dimension theorem proves its entire topological dimension is one.
\end{proof}

\begin{lemma}[Actual closed curve chains in schemes]\label{lem:fc-scheme-curve-chains}
\leanok
\lean{EnvelopeProperty.FcSchemeCurveChains.CurveStep, EnvelopeProperty.FcSchemeCurveChains.CurveChain, EnvelopeProperty.FcSchemeCurveChains.curveStep_map_closedImmersion, EnvelopeProperty.FcSchemeCurveChains.curveChain_map_closedImmersion, EnvelopeProperty.FcSchemeCurveChains.curveStep_of_affine, EnvelopeProperty.FcSchemeCurveChains.curveChain_of_affine, EnvelopeProperty.FcSchemeCurveChains.curveChain_of_common_affine, EnvelopeProperty.FcSchemeCurveChains.curveChain_of_integral, EnvelopeProperty.FcSchemeCurveChains.curveChain_of_common_irreducible_closed, EnvelopeProperty.FcSchemeCurveChains.curveChain_of_connected_finiteType}
\uses{lem:fc-affine-integral-curve-chains, lem:fc-scheme-curve-dimensions}
Let $F$ be algebraically closed and $X$ a connected finite-type $F$-scheme. Any two closed points
are joined by a finite chain of actual closed immersions $C_j\to X$, with each $C_j$ integral and
its ENTIRE global topological Krull dimension exactly one. No reducedness, irreducibility,
separatedness, cover, component, intersection-point, closure or dimension certificate is an input.
The empty chain is allowed when the endpoints coincide.
\end{lemma}

\begin{proof}
First suppose $X$ is integral. Choose actual affine neighbourhoods $U,V$ of the prescribed endpoints.
Their nonempty open intersection contains a closed point $z$, since the original finite-type scheme
is Jacobson. The actual affine-domain theorem, Lemma~\ref{lem:fc-affine-integral-curve-chains}, joins
both endpoint pairs inside their corresponding actual coordinate rings. Every prime quotient curve
in those chains closes by Lemma~\ref{lem:fc-scheme-curve-dimensions} to an actual integral CLOSED
curve in the ORIGINAL $X$ of global dimension one, with its entire original chart preimage unchanged.
Thus both chains have the exact original endpoints and glue through $z$.

For any actual irreducible closed subset $Z\subseteq X$, construct its integral reduced vanishing-ideal
subscheme. Prescribed original closed points in $Z$ lift to actual closed points of this subscheme.
The integral case joins them, and composition with its actual closed immersion into $X$ transports
all source curves, their global dimensions and both original endpoints back to $X$.

Finally, the actual finite-type structure morphism makes $X$ noetherian, so its irreducible components
are finite. Call a component good if it contains a closed point in the actual curve class of the first
endpoint. The union $U$ of good components and union $V$ of the others are closed and cover $X$.
They cannot intersect: any such component intersection is closed and contains a constructed actual
closed point. The preceding common-irreducible-closed-subset result joins it to the good-component
witness, so it also witnesses goodness of the remaining component, a contradiction. Thus $U$ is a
nonempty clopen. Connectedness forces $U=X$; the second endpoint lies on a good component and the
same actual common-component argument joins it to the first. Every step retains an actual source
scheme, actual CLOSED immersion into the ORIGINAL $X$, integrality, global dimension one, and
both original points in its actual image.
\end{proof}

\begin{lemma}[Chains of curves]\label{lem:fc-curve-chains}
\leanok
\lean{EnvelopeProperty.FcSchemeCurveChains.curveChain_of_connected_finiteType, EnvelopeProperty.FcProjectiveCurveChains.isProjective_closedImmersion_comp, EnvelopeProperty.FcProjectiveCurveChains.ProjectiveCurveStep, EnvelopeProperty.FcProjectiveCurveChains.ProjectiveCurveChain, EnvelopeProperty.FcProjectiveCurveChains.projectiveCurveChain_of_connected_finiteType}
\uses{lem:fc-scheme-curve-chains}
Let $F$ be algebraically closed and $X$ a connected $F$-scheme of finite type. Any two closed points $p,q$ of $X$ are
joined by a chain $C_1,\dots,C_l$ of irreducible closed curves, $p\in C_1$, $q\in C_l$, $C_j\cap C_{j+1}\ne\emptyset$. If $X$ is
projective, the $C_j$ may be taken projective. The empty chain is allowed when $p=q$.
\end{lemma}

\begin{proof}
Lemma~\ref{lem:fc-scheme-curve-chains} proves the full general statement, with actual closed integral
curve schemes and their entire global dimensions. If the actual structure morphism $f\colon X\to\Spec F$
is projective, choose its actual factorization $X\hookrightarrow\mathbb P(E)\to\Spec F$ by a quasi-coherent
module of finite type. For EVERY actual source curve $i\colon C_j\hookrightarrow X$, the composition
$C_j\hookrightarrow X\hookrightarrow\mathbb P(E)$ is an actual closed immersion with the same original
structure morphism $f\circ i$. It therefore constructs projectivity of that SAME curve $C_j$ over $F$,
retaining the actual original endpoints and intersection points. Apply this to every step of the checked
finite curve chain. No additional projective-curve or chain assumption is supplied.
\end{proof}

\begin{lemma}[Actual flat and local completed total sections]\label{lem:fc-proper-birational-flat-sections}
\leanok
\lean{EnvelopeProperty.FcProperBirationalFlatSections.flatSectionsEquiv, EnvelopeProperty.FcProperBirationalFlatSections.flatSectionsEquiv_apply, EnvelopeProperty.FcProperBirationalFlatSections.coefficientMap_flat_bijective, EnvelopeProperty.FcProperBirationalFlatSections.connectedSpace_of_globalSections_isLocalRing, EnvelopeProperty.FcProperBirationalFlatSections.connectedSpace_flatBaseChange, EnvelopeProperty.FcProperBirationalFlatSections.completedLocalRing, EnvelopeProperty.FcProperBirationalFlatSections.completedLocalBaseChange, EnvelopeProperty.FcProperBirationalFlatSections.completedLocalSectionsEquiv, EnvelopeProperty.FcProperBirationalFlatSections.connectedSpace_completedLocalBaseChange, EnvelopeProperty.FcProperBirationalFlatSections.completedLocal_idempotents_trivial}
\uses{scheme:actual-reduced-proper-birational-global-sections, scheme:actual-original-flat-global-section-tensor}
Let $A$ be an integrally closed domain, $f\colon Y\to\Spec A$ an actual proper birational map with
reduced source, and $S$ an actual flat $A$-algebra. The ORIGINAL whole coefficient map gives a constructed
$S$-algebra equivalence $S\simeq\Gamma(Y_S,\mathcal O)$. If $S$ is local, the actual TOTAL scheme $Y_S$
is connected. If $A$ is noetherian, for every actual prime $P$ take the actual maximal-adic completion
$\widehat{A_P}$: its flatness and locality are derived, and the entire original total section ring
is $\widehat{A_P}$ with every whole global idempotent equal to zero or one.
\end{lemma}

\begin{proof}
The actual proper birational normal-base theorem constructs the entire coefficient equivalence
$A\simeq\Gamma(Y,\mathcal O)$. Properness supplies the genuine QCQS hypotheses, so the actual
flat section pushout identifies $\Gamma(Y_S,\mathcal O)$ with $S\otimes_A\Gamma(Y,\mathcal O)$.
Tensor the constructed entire original coefficient equivalence and compose with the tensor-unit
equivalence. Its whole $S$-coefficient map is exactly the actual original pullback coefficient map.
If $S$ is local, this makes the whole global-section ring local. Its actual global idempotents
are zero or one by the local-ring unit alternative for $e$ and $1-e$; nontrivial global sections
give actual scheme nonemptiness. The genuine global-idempotent/clopen correspondence therefore
proves connectedness. At any actual prime, localization preserves noetherianity and is flat;
the actual maximal-adic completion is flat and local. Compose the two actual flat structures
and apply the preceding total-scheme construction. No normality or domain property of that
actual completion is assumed.

This proves connectedness of the COMPLETED TOTAL scheme, not of its closed fibre. The remaining
formal-functions step must algebraize actual compatible infinitesimal fibre idempotents into
whole global idempotents of this completed total scheme. No such surjectivity, fibre connectedness,
or equality of the nonreduced fibre's entire section ring with its coefficient field is asserted here.
\end{proof}

\begin{lemma}[Actual flat local lifts of every prescribed field point]\label{lem:fc-flat-local-field-lift}
\leanok
\lean{EnvelopeProperty.FcFlatLocalFieldLift.tensorEvaluation, EnvelopeProperty.FcFlatLocalFieldLift.tensorEvaluation_surjective, EnvelopeProperty.FcFlatLocalFieldLift.tensorEvaluation_comp, EnvelopeProperty.FcFlatLocalFieldLift.isNoetherianRing_tensor, EnvelopeProperty.FcFlatLocalFieldLift.tensorEvaluation_ker_isMaximal, EnvelopeProperty.FcFlatLocalFieldLift.exists_noetherian_flat_local_lift, EnvelopeProperty.FcFlatLocalFieldLift.completionEvaluation, EnvelopeProperty.FcFlatLocalFieldLift.completionEvaluation_algebraMap, EnvelopeProperty.FcFlatLocalFieldLift.completionEvaluation_surjective, EnvelopeProperty.FcFlatLocalFieldLift.completionEvaluation_ker, EnvelopeProperty.FcFlatLocalFieldLift.exists_completed_flat_local_lift}
\uses{lem:fc-proper-birational-flat-sections}
Let $F$ be a field, $A$ a finite-type $F$-algebra and $w\colon A\to K$ any actual $F$-algebra point
in any field extension $K/F$. One can construct a noetherian local flat $A$-algebra $R$ and a surjection
$\beta\colon R\to K$ with kernel EXACTLY $\mathfrak m_R$ and entire composite $A\to R\to K$ equal to $w$.
Its actual maximal-adic completion $B=\widehat R$ is flat over $A$, local and complete, with a constructed
surjection $\gamma\colon B\to K$ whose kernel is EXACTLY $\mathfrak m_B$ and whose entire original
$A$-coefficient composite is again $w$. No finite or separable extension, scalar lift, local ring,
completion or residue-field identification is supplied as input.
\end{lemma}

\begin{proof}
Form the ACTUAL tensor algebra $T=A\otimes_F K$ and its WHOLE evaluation
$\theta(a\otimes k)=w(a)k$. It is surjective since $\theta(1\otimes k)=k$ and retains the entire
original coefficient map. Thus $P=\ker\theta$ is an actual maximal ideal. The tensor ring is flat
over $A$, and noetherian because, through the actual tensor interchange isomorphism, it is finite
type over the prescribed field $K$; no finiteness of $K/F$ is required. Take the actual localization
$R=T_P$. Its whole evaluation extends uniquely since every denominator outside $P$ has nonzero
field image. This extension is still surjective, so its kernel is the unique actual maximal ideal of $R$.
Localization preserves noetherianity and flatness, and its entire original $A$-coefficient map remains $w$.

Complete the ACTUAL maximal ideal of that ACTUAL noetherian local $R$. Genuine adic-completion
flatness, locality and completeness construct the required complete local flat $A$-algebra $B$.
The original evaluation descends to $R/\mathfrak m_R=K$; compose with the ACTUAL order-one
completion quotient $B\to R/\mathfrak m_R$. Every original local coefficient retains its whole field
value. In particular this completed evaluation is surjective, so its kernel is the actual unique maximal
ideal of $B$. All whole original coefficient triangles are proved. Thus every prescribed geometric
field fibre has the exact coefficient field of a constructed noetherian local scalar lift and its completion.
The remaining proper formal-functions/idempotent algebraization step is not assumed or asserted.
\end{proof}

\begin{lemma}[Actual kernels commute with finite-module completion]\label{lem:fc-adic-completion-kernel}
\leanok
Let $R$ be noetherian, $I$ an actual ideal and $f\colon M\to N$ an actual linear map of finite
$R$-modules. The ACTUAL completed kernel inclusion constructs an $\widehat R$-linear equivalence
$\widehat{\ker f}\simeq\ker(\widehat f)$, retaining the WHOLE completed inclusion and every original
coefficient section. More generally completion preserves exactness of $L\to M\to N$ when $M,N$
are finite, with no injectivity or surjectivity hypothesis on either actual differential.
\end{lemma}

\begin{proof}
Factor $f$ through its ACTUAL range. Apply genuine Artin--Rees short-exact completion to
$0\to\ker f\to M\to\operatorname{range} f\to0$. Completion of the actual range inclusion
into the finite module $N$ is injective, again by Artin--Rees. Thus the completed original
kernel inclusion has precisely the actual kernel of the completed original map as image,
and is injective. Its actual codomain restriction supplies the equivalence, whose underlying
whole map and original coefficient sections are proved to be the literal completed inclusion.
For an arbitrary exact complex, its original left differential surjects onto the actual right
kernel; completion preserves that surjection. Combine it with the constructed kernel
equivalence to prove exactness without extra hypotheses on the differentials.

This is genuine finite-module algebra. Raw affine \v Cech terms of a proper scheme need not
be finite over the base. An actual proper coherent-finiteness/finite-complex replacement and
the formal-H0 comparison must still be derived geometrically; neither is an assumption or a
conclusion of this lemma. In particular Theorem~\ref{thm:fc-zariski} remains open.
\end{proof}

\begin{lemma}[Actual coefficient extension of Rees algebras]\label{lem:fc-rees-ideal-base-change}
\leanok
\lean{EnvelopeProperty.FcReesIdealBaseChange.reesMap, EnvelopeProperty.FcReesIdealBaseChange.baseChangeMap, EnvelopeProperty.FcReesIdealBaseChange.baseChangeMap_tmul_val, EnvelopeProperty.FcReesIdealBaseChange.baseChangeMap_tmul_monomial_val, EnvelopeProperty.FcReesIdealBaseChange.baseChangeMap_surjective, EnvelopeProperty.FcReesIdealBaseChange.baseChangeMap_comp_includeRight, EnvelopeProperty.FcReesIdealBaseChange.baseChangeMap_comp_algebraMap, EnvelopeProperty.FcReesIdealBaseChange.spec_baseChangeMap_isClosedImmersion}
For ANY actual ring algebra $R\to A$ and actual ideal $I\subseteq R$, the genuine coefficient
polynomial map constructs an actual surjection
$A\otimes_R R[It]\twoheadrightarrow A[(IA)t]$ and hence an actual closed immersion of their
coefficient affine schemes. No flatness, injectivity, noetherianity or pre-supplied Rees
comparison is required. The WHOLE $A$ coefficient and original Rees triangles commute;
every homogeneous monomial $a\otimes rt^n$ maps to $a\overline r t^n$ with its actual degree.
\end{lemma}

\begin{proof}
The whole original coefficient polynomial map carries each actual $I^n$ coefficient into
$(IA)^n$, by extension of ideal powers, so it genuinely restricts to the two actual Rees
subalgebras. The actual tensor universal property constructs the displayed map, retaining
all coefficient maps and homogeneous monomials. The extended Rees algebra is generated over
$A$ by the actual degree-one monomials of $IA$. The elements of $IA$ are generated over $A$
by the images of actual elements of $I$, and each corresponding degree-one monomial is the
image of its actual original Rees tensor. Thus the whole polynomial image contains the whole
extended Rees algebra. Its actual affine spectrum map is therefore a closed immersion.
This gives genuine affine pieces toward the coherent Rees family in formal functions; global
Rees-sheaf gluing and the proper formal-H0 comparison remain to be proved.
\end{proof}

\begin{lemma}[Actual proper schematic Rees family]\label{lem:fc-rees-formal-family}
\leanok
\lean{EnvelopeProperty.FcReesFormalFamily.reesBase, EnvelopeProperty.FcReesFormalFamily.polynomialToRees, EnvelopeProperty.FcReesFormalFamily.polynomialToRees_comp_base, EnvelopeProperty.FcReesFormalFamily.ambient, EnvelopeProperty.FcReesFormalFamily.ambientProjection, EnvelopeProperty.FcReesFormalFamily.polynomialSource, EnvelopeProperty.FcReesFormalFamily.imageMap, EnvelopeProperty.FcReesFormalFamily.family, EnvelopeProperty.FcReesFormalFamily.familyImmersion, EnvelopeProperty.FcReesFormalFamily.familyProjection, EnvelopeProperty.FcReesFormalFamily.imageMap_isAffineHom, EnvelopeProperty.FcReesFormalFamily.familyProjection_isProper, EnvelopeProperty.FcReesFormalFamily.polynomialSource_isPullback, EnvelopeProperty.FcReesFormalFamily.family_isNoetherian, EnvelopeProperty.FcReesFormalFamily.family_isSeparated, EnvelopeProperty.FcReesFormalFamily.polynomialToFamily, EnvelopeProperty.FcReesFormalFamily.polynomialToFamily_comp_projection}
\uses{lem:fc-rees-ideal-base-change}
For any actual $f\colon X\to\operatorname{Spec}R$ and actual ideal $I$, construct the actual
Rees coefficient extension $W=X\times_R\operatorname{Spec}R[It]$ and the actual polynomial
map $X\times_R\mathbb A^1_R\to W$. Its GENUINE scheme-theoretic image $Z\hookrightarrow W$
has a proper projection to $\operatorname{Spec}R[It]$ whenever $f$ is proper. If $R$ is
noetherian, $Z$ is noetherian; it is separated as an absolute scheme. The whole polynomial
coefficient triangle and the original polynomial-base-change square are retained.
No family, image comparison, properness certificate or cohomological assertion is input.
\end{lemma}

\begin{proof}
Use the ACTUAL inclusion $R[It]\hookrightarrow R[t]$ and the literal iterated pullback along
its actual affine spectrum map. Pullback pasting, together with the proved ENTIRE original
$R$ coefficient triangle, identifies its source with the original $X\times_R\mathbb A^1_R$.
The induced map to $W$ is affine by genuine base change of an actual affine morphism.
Take its ACTUAL kernel ideal sheaf and associated subscheme. This is the literal schematic
image, with all nilpotents retained, and its inclusion is a genuine closed immersion.
Its projection is proper as a closed immersion followed by the ACTUAL proper base change
of $f$. Over noetherian $R$, the actual Rees algebra is noetherian, and local finite type and
quasi-compactness give noetherianity of the actual image scheme. Properness over an affine
scheme gives absolute separatedness. The original polynomial map factors through its actual
image by the genuine image universal construction, retaining the WHOLE coefficient triangle.

The actual affine chart identification with $A[(IA)t]$, its graded cohomology comparison and
formal-H0 remain to be derived. This lemma constructs the family and its genuine geometric
source properties; it does not assert the missing formal-functions theorem.
\end{proof}

\begin{lemma}[Actual affine polynomial image is the extended Rees scheme]\label{lem:fc-rees-formal-affine}
\leanok
\lean{EnvelopeProperty.FcReesFormalAffine.polynomialTensorEquiv, EnvelopeProperty.FcReesFormalAffine.polynomialTensorEquiv_tmul, EnvelopeProperty.FcReesFormalAffine.tensorInclusion, EnvelopeProperty.FcReesFormalAffine.polynomialTensorEquiv_comp_inclusion, EnvelopeProperty.FcReesFormalAffine.tensorInclusion_ker, EnvelopeProperty.FcReesFormalAffine.affineImageEquiv, EnvelopeProperty.FcReesFormalAffine.affineImageEquiv_mk, EnvelopeProperty.FcReesFormalAffine.specMap_ker_eq_of_ringKer_eq, EnvelopeProperty.FcReesFormalAffine.affineImageSchemeIso, EnvelopeProperty.FcReesFormalAffine.affineImageSchemeIso_comp_immersion}
\uses{lem:fc-rees-ideal-base-change, lem:fc-rees-formal-family}
For any $R\to A$ and actual ideal $I$, construct the actual polynomial tensor isomorphism
$A\otimes_R R[t]\simeq A[t]$. Under this WHOLE isomorphism the genuine tensor extension of
$R[It]\hookrightarrow R[t]$ factors through the actual surjection onto $A[(IA)t]$.
Consequently its ACTUAL affine schematic-image quotient is $A[(IA)t]$ as an $A$-algebra,
and its genuine SCHEME-theoretic image is isomorphic to $\operatorname{Spec}A[(IA)t]$,
retaining the ENTIRE closed immersion into $\operatorname{Spec}(A\otimes_R R[It])$.
No flatness, reducedness or pre-supplied image comparison is required.
\end{lemma}

\begin{proof}
Construct the polynomial tensor isomorphism from the genuine one-variable multivariate
polynomial tensor equivalence and the two actual one-variable polynomial equivalences.
Polynomial induction proves its whole tensor formula $a\otimes p\mapsto a\overline p$.
Compare the actual tensor inclusion and actual extended Rees surjection on every pure
tensor to prove equality of the ENTIRE algebra maps after transport. Injectivity of the
polynomial isomorphism and actual Rees subalgebra inclusion gives equality of the two actual
kernels. The genuine quotient first-isomorphism theorem constructs the whole $A$-algebra
image comparison, preserving every original tensor section.

Equality of actual affine ring kernels gives equality of the genuine scheme-image ideal
sheaves: naturality of the actual $\Gamma\operatorname{Spec}$ isomorphisms transports every
whole affine coefficient kernel. The actual extended Rees map is a closed immersion by
its proved surjectivity. The universal property of closed immersions and equality of the
actual kernels therefore construct a genuine isomorphism of schematic images; its ENTIRE
ambient closed-immersion triangle is proved. Transport to the literal iterated global family
and subsequent graded cohomological/formal-H0 comparison remain to be derived.
\end{proof}

\begin{lemma}[Literal affine-base schematic Rees family]\label{lem:fc-rees-formal-affine-scheme}
\leanok
\lean{EnvelopeProperty.FcReesFormalAffineScheme.affineBase, EnvelopeProperty.FcReesFormalAffineScheme.ambientIso, EnvelopeProperty.FcReesFormalAffineScheme.polynomialSourceIso, EnvelopeProperty.FcReesFormalAffineScheme.polynomialSourceIso_hom_fst, EnvelopeProperty.FcReesFormalAffineScheme.polynomialSourceIso_hom_snd, EnvelopeProperty.FcReesFormalAffineScheme.spec_tensorInclusion_fst, EnvelopeProperty.FcReesFormalAffineScheme.spec_tensorInclusion_snd, EnvelopeProperty.FcReesFormalAffineScheme.imageMap_comp_ambientIso, EnvelopeProperty.FcReesFormalAffineScheme.imageIsoOfArrowIso, EnvelopeProperty.FcReesFormalAffineScheme.imageIsoOfArrowIso_comp, EnvelopeProperty.FcReesFormalAffineScheme.familyIso, EnvelopeProperty.FcReesFormalAffineScheme.familyIso_comp_immersion, EnvelopeProperty.FcReesFormalAffineScheme.familyIso_comp_projection}
\uses{lem:fc-rees-formal-family, lem:fc-rees-formal-affine}
For any actual $R$-algebra $A$ and actual ideal $I$, the LITERAL global schematic Rees family
constructed from $\operatorname{Spec}A\to\operatorname{Spec}R$ is genuinely isomorphic to
$\operatorname{Spec}A[(IA)t]$. The isomorphism retains BOTH the ENTIRE closed immersion into
the actual Rees coefficient base extension and the ENTIRE original Rees coefficient projection.
No flatness, reducedness, noetherianity or supplied image comparison is required.
\end{lemma}

\begin{proof}
Actual affine pullback tensor isomorphisms identify the literal ambient scheme with
$\operatorname{Spec}(A\otimes_R R[It])$. Actual pullback pasting identifies the literal polynomial
source with $\operatorname{Spec}(A\otimes_R R[t])$, retaining BOTH whole coefficient projections.
The actual polynomial-to-Rees global morphism is proved equal to the genuine tensor inclusion
under these two whole isomorphisms, by the actual pullback universal property and the two
entire coefficient diagrams. Isomorphisms of actual source/target arrows induce isomorphisms
of genuine scheme-theoretic images: compare the actual closed immersion kernels and apply
the genuine closed-immersion lifting universal property. Combine with the checked affine
image identification from Lemma~\ref{lem:fc-rees-formal-affine}. The ENTIRE ambient inclusion
and Rees coefficient projection equalities are then proved as actual scheme morphisms.
General open restriction charts and subsequent graded cohomological/formal-H0 comparison
remain to be derived; no formal-functions conclusion is asserted here.
\end{proof}

\begin{lemma}[Actual open restriction and affine charts of the schematic Rees family]\label{lem:fc-rees-formal-restriction}
\leanok
\lean{EnvelopeProperty.FcReesFormalRestriction.ker_eq_comap_of_open_pullback, EnvelopeProperty.FcReesFormalRestriction.imageRestrictionMap, EnvelopeProperty.FcReesFormalRestriction.isPullback_imageRestriction, EnvelopeProperty.FcReesFormalRestriction.ambientMap, EnvelopeProperty.FcReesFormalRestriction.ambientMap_projection, EnvelopeProperty.FcReesFormalRestriction.ambientMap_isOpenImmersion, EnvelopeProperty.FcReesFormalRestriction.polynomialSourceMap, EnvelopeProperty.FcReesFormalRestriction.polynomialSourceMap_imageMap, EnvelopeProperty.FcReesFormalRestriction.polynomialSourceMap_snd, EnvelopeProperty.FcReesFormalRestriction.polynomialSourceMap_isPullback, EnvelopeProperty.FcReesFormalRestriction.familyMap, EnvelopeProperty.FcReesFormalRestriction.familyMap_isPullback, EnvelopeProperty.FcReesFormalRestriction.familyMap_isOpenImmersion, EnvelopeProperty.FcReesFormalRestriction.ambientMap_fst, EnvelopeProperty.FcReesFormalRestriction.ambientMap_isPullback, EnvelopeProperty.FcReesFormalRestriction.familyMap_isPullback_over_original, EnvelopeProperty.FcReesFormalRestriction.familyMap_projection, EnvelopeProperty.FcReesFormalRestriction.affineFamilyIso_comp_original, EnvelopeProperty.FcReesFormalRestriction.chartCoefficient, EnvelopeProperty.FcReesFormalRestriction.chartAlgebra, EnvelopeProperty.FcReesFormalRestriction.affineBase_eq_chart, EnvelopeProperty.FcReesFormalRestriction.chartFamilyIso, EnvelopeProperty.FcReesFormalRestriction.chartFamilyIso_comp_projection, EnvelopeProperty.FcReesFormalRestriction.chartMap, EnvelopeProperty.FcReesFormalRestriction.chartMap_isOpenImmersion, EnvelopeProperty.FcReesFormalRestriction.chartMap_projection, EnvelopeProperty.FcReesFormalRestriction.chartFamilyIso_comp_original, EnvelopeProperty.FcReesFormalRestriction.chartMap_isPullback_over_original, EnvelopeProperty.FcReesFormalRestriction.chartMap_opensRange, EnvelopeProperty.FcReesFormalRestriction.affineChartCover}
\uses{lem:fc-rees-formal-family, lem:fc-rees-formal-affine-scheme}
The actual schematic Rees family commutes with every original open immersion as a genuine
pullback over the original scheme. For an actual affine open $U$, its whole inverse image
is the actual affine scheme $\operatorname{Spec}\Gamma(U)[(I\Gamma(U))t]$. These actual
charts form an affine open cover of the whole family. Both original-scheme and Rees
coefficient morphisms are retained on all functions. No flatness, reducedness,
noetherianity or separately supplied chart comparison is required.
\end{lemma}

\begin{proof}
The actual kernel ideal sheaf of a quasi-compact morphism commutes with genuine open
pullback. The actual closed-immersion universal property therefore constructs the map
of schematic images and proves that its entire closed-immersion square is a pullback.
For the polynomial-to-Rees morphism, literal iterated pullback maps construct both
ambient and polynomial-source restriction maps; cancellation and pasting of genuine
pullbacks prove the required source square. Passing to the actual schematic image
then gives an open immersion of the constructed families, with both whole projections.
For an original affine open $U$, construct its coefficient homomorphism directly by
full faithfulness of $\operatorname{Spec}$ applied to the original chart morphism.
Lemma~\ref{lem:fc-rees-formal-affine-scheme} identifies the resulting actual affine
family with its actual extended-ideal Rees spectrum. Its canonical open inclusion is
proved to be the full inverse image of $U$, retaining both entire coefficient triangles.
The original affine cover therefore yields the actual affine Rees cover of the entire
family. Graded cohomological comparison and proper formal-H0 algebraization remain
unproved here; no connected-fibre conclusion is asserted.
\end{proof}

\begin{lemma}[Actual Rees chart sections and degree-preserving restriction]\label{lem:fc-rees-formal-chart-sheaf}
\leanok
\lean{EnvelopeProperty.FcReesFormalChartSheaf.projection, EnvelopeProperty.FcReesFormalChartSheaf.projection_isAffineHom, EnvelopeProperty.FcReesFormalChartSheaf.chartSectionsIso, EnvelopeProperty.FcReesFormalChartSheaf.chartCoefficient_naturality, EnvelopeProperty.FcReesFormalChartSheaf.reesRestriction, EnvelopeProperty.FcReesFormalChartSheaf.reesRestriction_val, EnvelopeProperty.FcReesFormalChartSheaf.reesRestriction_coeff, EnvelopeProperty.FcReesFormalChartSheaf.reesCoefficient, EnvelopeProperty.FcReesFormalChartSheaf.reesRestriction_comp_algebraMap, EnvelopeProperty.FcReesFormalChartSheaf.reesRestriction_comp_reesCoefficient, EnvelopeProperty.FcReesFormalChartSheaf.chartMap_original, EnvelopeProperty.FcReesFormalChartSheaf.chartMap_rees, EnvelopeProperty.FcReesFormalChartSheaf.spec_reesRestriction_comp_algebraMap, EnvelopeProperty.FcReesFormalChartSheaf.chartMap_naturality, EnvelopeProperty.FcReesFormalChartSheaf.chartMap_preimage, EnvelopeProperty.FcReesFormalChartSheaf.chartSectionsIso_hom, EnvelopeProperty.FcReesFormalChartSheaf.chartSectionsIso_naturality}
\uses{lem:fc-rees-formal-restriction, lem:fc-rees-ideal-base-change}
The actual original-scheme projection $q$ of the schematic Rees family is affine.
For every original affine open $U$, there is a canonical isomorphism
$\Gamma(q^{-1}U,\mathcal O)\simeq\Gamma(U)[(I\Gamma(U))t]$.
The actual sheaf restriction for $U\subseteq V$ becomes coefficientwise polynomial
restriction under these isomorphisms, preserving each actual degree and both entire
original coefficient homomorphisms. No separately supplied compatibility, flatness,
reducedness or noetherianity is required.
\end{lemma}

\begin{proof}
Every original affine-open inverse image is affine by the actual chart pullback and
range theorem, proving $q$ affine. The actual open-immersion section isomorphism and
the genuine $\Gamma\operatorname{Spec}$ isomorphism give the full chart-section
comparison. Full faithfulness of $\operatorname{Spec}$ and the original affine-open
restriction identity prove compatibility of the entire original chart coefficients.
Coefficientwise polynomial restriction then maps every extended ideal power into
the corresponding actual ideal power, defining the true Rees ring restriction.
The full degree-zero and original Rees coefficient triangles commute. The original
closed immersion into the ambient coefficient pullback is a monomorphism, so these
two entire triangles prove that its spectrum is the actual original global chart
restriction. Naturality of actual scheme section maps and of
$\Gamma\operatorname{Spec}$ now prove the entire sheaf-restriction identity.
The subsequent graded cohomological comparison and formal-H0 algebraization remain
to be proved; this lemma asserts no connected-fibre theorem.
\end{proof}

\begin{lemma}[Actual ideal-power decomposition of Rees chart sections]\label{lem:fc-rees-power-decomposition}
\leanok
\lean{EnvelopeProperty.FcReesPowerDecomposition.coefficients, EnvelopeProperty.FcReesPowerDecomposition.coefficients_apply, EnvelopeProperty.FcReesPowerDecomposition.coefficientLinear, EnvelopeProperty.FcReesPowerDecomposition.coefficientLinear_injective, EnvelopeProperty.FcReesPowerDecomposition.coefficients_monomial, EnvelopeProperty.FcReesPowerDecomposition.coefficientLinear_surjective, EnvelopeProperty.FcReesPowerDecomposition.powerEquiv, EnvelopeProperty.FcReesPowerDecomposition.finitePiDirectSumEquiv, EnvelopeProperty.FcReesPowerDecomposition.finitePiDirectSumEquiv_apply, EnvelopeProperty.FcReesPowerDecomposition.finiteReesPowerEquiv, EnvelopeProperty.FcReesPowerDecomposition.finiteReesPowerEquiv_apply, EnvelopeProperty.FcReesPowerDecomposition.chartReesRestriction, EnvelopeProperty.FcReesPowerDecomposition.chartPowerRestriction, EnvelopeProperty.FcReesPowerDecomposition.chartPowerEquiv_naturality}
\uses{lem:fc-rees-formal-chart-sheaf}
The genuine Rees algebra is linearly the direct sum of its actual ideal powers,
with every original polynomial coefficient retained in its original degree.
For any finite family of actual chart algebras over the original base, the product
of their Rees algebras is linearly the direct sum of the products of their actual
ideal powers. Under these genuine decompositions, actual original affine-chart
restriction is the direct sum of the actual original ideal-power restrictions,
linear over the original coefficient ring. No noetherianity, flatness or assumed
graded-comparison certificate is required.
\end{lemma}

\begin{proof}
Send an actual Rees polynomial to its finitely supported family of actual
coefficients in the corresponding ideal powers. This map is linear and injective
by polynomial extensionality. A homogeneous monomial maps exactly to the direct-sum
inclusion of its coefficient. Direct-sum induction proves surjectivity and hence
constructs the genuine linear equivalence. Finite products commute with direct sums
by replacing the finite product with a finite direct sum and exchanging the two
indices; the resulting equivalence is proved to retain every index and degree.
Apply this comparison to the actual ideal powers of all chart algebras. Actual
coefficientwise Rees restriction respects the derived original-base algebra
structures, and its individual degree restrictions are linear over that same base.
Equality on every actual coefficient proves the full direct-sum restriction square.
The differential and cohomological comparison of the actual ordered Cech complexes
still remain to be derived.
\end{proof}

\begin{theorem}[Connected fibres {\cite[Theorem~6.1]{FHX}}]\label{thm:fc-zariski}
\lean{EnvelopeProperty.ProjectiveGeometricConnectedness.geometricallyConnected,EnvelopeProperty.ProjSubvariety.BlowupCenter.fcNativeRees_geometricallyConnected}
\leanok
\uses{def:pg-blowup, lem:fc-proper-birational-flat-sections, lem:fc-flat-local-field-lift}
If $T'$ is normal, every scheme-theoretic fibre $B_\xi$ of $B\to T'$ is geometrically connected.
More generally, a reduced projective birational scheme over the spectrum of a finite-type
normal domain has geometrically connected fibres.
\end{theorem}
\begin{proof}
For each original field-valued point, construct the complete flat local model with that
same residue field. Flat base change of the original proper birational coefficient
bijection gives the whole global section ring of this model as its local base ring.
Actual projective formal H0 identifies it with all compatible original infinitesimal
sections. Unique lifting through the genuine nilpotent thickenings identifies the
idempotents of the whole closed fibre with those of that local ring. They are zero and
one, and nontriviality is retained, so the whole nonreduced closed fibre is connected.
The constructed scheme isomorphism identifies it with the literal original field-valued
fibre. Since this holds for every such field-valued point, standard geometric connectedness
follows. Normality of the completion and equality of closed-fibre global sections with its
residue field are not used.
\end{proof}

\begin{theorem}[Maximum principle on curves {\cite[Proposition~3.5]{GR25}, \cite[Lemma~5.2]{FHX}}]\label{thm:fc-curve-mp}
\leanok
\uses{thm:fc-connected-finite-type-curve-maximum}
Let $F$ be a complete algebraically closed ultrametric normed field and $D$ a connected finite-type $F$-scheme of global dimension at most one. A native locally psh function on $D^\an$ attaining its supremum is constant. This strengthens the smooth connected affine-curve case and includes the trivial base norm.
\end{theorem}

\begin{proof}
Take the reduced singleton closure of a maximizing support and lift the entire original native point to that integral closed curve. The integral-curve maximum theorem makes the pullback constant. A genuine closed rational section gives a closed original support with the same maximum. For any other point, take its integral reduced support closure and another closed rational section. Actual finite-type curve chains with closed intermediates compare the two rational anchors; the integral-curve maximum theorem on the second support closure then compares the arbitrary point with its anchor. Every lift, closed subscheme and rational section is constructed from the original scheme.
\end{proof}

\begin{theorem}[Global maximum principle]\label{thm:fc-global-mp}
\lean{EnvelopeProperty.Berkovich.RelativeSchemeValuedPoint.IsLocallyPSHOn.eq_of_isMaxOn_univ_geometricallyConnected_proper}
\leanok
\uses{thm:fc-geometrically-connected-proper-native-maximum}
Let $F$ be a complete ultrametric normed field, $X$ a geometrically connected proper $F$-scheme, and $u$ native locally psh on $X^\an$. If $u$ attains its maximum $R$, then $u\equiv R$. Properness strengthens the projective case used in the paper.
\end{theorem}

\begin{proof}
Pass to the constructed completed algebraic closure of $F$. Actual scalar restriction is surjective on whole native points, preserves local psh, and lifts the original maximizing point. Properness and geometric connectedness give a connected proper scalar-changed scheme. Over the algebraically closed field, each arbitrary support is compared with an actual rational section of its integral reduced singleton closure: the original coefficient map of that proper integral scheme is bijective, and genuine field base change gives its geometric connectedness. Closed-support curve chains then propagate the attained maximum, and the two rational comparisons recover every original native value. Finally descend through the surjective scalar restriction. No rational-lift, curve-chain or geometric-connectedness certificate is supplied beyond the theorem's original geometric hypothesis.
\end{proof}

\begin{proof}[Proof of Theorem~\ref{thm:gl-fibre-constancy}]
\uses{lem:fc-fibre-model, lem:fc-fibre-psh, thm:fc-zariski, thm:fc-global-mp}
Let $t\in T'^\an$ with support $\xi$, and choose $(m,s)$ with $|s|_t\ne0$. By Lemma~\ref{lem:fc-fibre-psh} it suffices to show
that $u_t$ is constant. $Z_t=B_\xi\otimes\mathcal H(t)$ is projective and geometrically connected
(Theorem~\ref{thm:fc-zariski}), and $Z_t^\an$ is compact. If $u_t\equiv-\infty$ there is nothing to prove; otherwise the
usc function $u_t$ attains a finite maximum, and Theorem~\ref{thm:fc-global-mp} applies.
\end{proof}

\begin{remark}[A countability-free analytic route]\label{rmk:fc-directed-analytic-route}
\uses{thm:fc-analytic-constancy-completed-geometrically-connected,cor:fc-analytic-constancy-original-geometrically-connected-rees}
The checked native analytic reduction also has a directed-net proof. The original global
blowup psh potential supplies a decreasing dyadic finite Fubini--Study net, whose actual
Rees-chart restrictions are convex PL nets in regular units. Their whole native fibre
pullback satisfies the directed maximum principle, with compactness and an attained
maximum derived internally. Thus standard geometric connectedness of the literal original
completed fibre suffices for whole original analytic fibre constancy, without countability
of the base field. This does not identify arbitrary directed presentations with the
sequence-based local psh definition in Lemma~\ref{lem:fc-fibre-psh}, and it does not yet prove
Theorem~\ref{thm:fc-zariski}.
\end{remark}

\subsection{Potential theory on the Berkovich line}\label{subsec:fc-line}

\begin{convention}\label{conv:fc-line}
\notready
$F$ is complete, algebraically closed and nontrivially valued, and $\A^1=\A^1_F$. For $a\in F$ and $r\ge0$, $\zeta_{a,r}$ is
the Gauss point $|\sum c_n(T-a)^n|=\max|c_n|r^n$. $D(a,r)=\{|T-a|\le r\}$, $D^-(a,r)=\{|T-a|<r\}$.
\end{convention}

\begin{definition}[Standard affinoids]\label{def:fc-standard-affinoid}
\lean{EnvelopeProperty.Berkovich.RelativeSpectrum.FCStandardAffinoid, EnvelopeProperty.Berkovich.RelativeSpectrum.FCStandardAffinoid.boundary, EnvelopeProperty.Berkovich.RelativeSpectrum.FCStandardAffinoid.radialSkeleton}
\leanok
\uses{conv:fc-line}
A \emph{standard affinoid} is $U=D(a_0,r_0)\setminus\bigcup_{j=1}^mD^-(a_j,r_j)$ with $0<r_j\le r_0$, $|a_j-a_0|\le r_0$, and the
$D^-(a_j,r_j)$ pairwise disjoint. Its boundary is $\partial U=\{\zeta_{a_j,r_j}\}_{0\le j\le m}$, and its skeleton $\Sigma(U)$ is
the union of the segments $[\zeta_{a_j,r_j},\zeta_{a_0,r_0}]$.
\end{definition}

\begin{lemma}[Topology of the line {\cite[Appendix~A]{GR25}}]\label{lem:fc-line-topology}
\lean{EnvelopeProperty.Berkovich.RelativeSpectrum.fc_gauss_or_nested_discs, EnvelopeProperty.Berkovich.RelativeSpectrum.fc_continuous_centeredGauss_prod, EnvelopeProperty.Berkovich.RelativeSpectrum.FCStandardAffinoid.isCompact_carrier, EnvelopeProperty.Berkovich.RelativeSpectrum.FCStandardAffinoid.isConnected_carrier, EnvelopeProperty.Berkovich.RelativeSpectrum.FCStandardAffinoid.interior_carrier, EnvelopeProperty.Berkovich.RelativeSpectrum.FCStandardAffinoid.exists_boundary_mem_closure_component, EnvelopeProperty.Berkovich.RelativeSpectrum.FCStandardAffinoid.radialRetraction, EnvelopeProperty.Berkovich.RelativeSpectrum.FCStandardAffinoid.radialRetraction_eq_self, EnvelopeProperty.Berkovich.RelativeSpectrum.fc_nhds_hasBasis_standardAffinoid, EnvelopeProperty.Berkovich.RelativeSpectrum.fc_exists_annulus_subset_nhds, EnvelopeProperty.Berkovich.RelativeSpectrum.fc_denseRange_rigid_norm, EnvelopeProperty.Berkovich.RelativeSpectrum.fc_isConnected_line}
\leanok
\uses{def:fc-standard-affinoid, thm:fc-nested-disc-classification, lem:fc-linear-topology, lem:fc-standard-inflation, lem:fc-standard-boundary, lem:fc-standard-component-boundary, thm:fc-radial-retraction, thm:fc-residue-end-annulus, cor:fc-line-paths, thm:fc-rigid-density, thm:fc-standard-affinoid-basis}
\begin{enumerate}
\item Points are $\zeta_{a,r}$ ($r\ge0$) or limits of Gauss points of nested classical discs with empty classical intersection; $(a,r)\mapsto\zeta_{a,r}$ is
continuous.
\item A standard affinoid $U$ is compact and connected, and its topological interior is $U\setminus\partial U$. Each
connected component of $U\setminus\partial U$ has a boundary point in its closure. There is a retraction
$\tau_U\colon U\to\Sigma(U)$.
\item Every point has a basis of neighbourhoods $\operatorname{int}(U)$, $U$ standard. At $\zeta_{a,r}$ ($r>0$) these
neighbourhoods contain, in each residue class $D^-(b,r)$, an annulus $\{\rho<|T-b|<r\}$.
\item Rigid points are dense, and $(\A^1)^\an$ is connected.
\end{enumerate}
\end{lemma}

\begin{proof}
(1) The infimum of the linear values is either attained, giving a Gauss point, or is approached by a strictly decreasing sequence. The corresponding discs are nested with empty classical intersection and their actual Gauss points converge by continuous inflation (Theorem~\ref{thm:fc-nested-disc-classification}). Joint continuity is Lemma~\ref{lem:fc-linear-topology}.

(2) Inflation gives paths to the outer Gauss point, compactness follows from the closed coordinate inequalities, and the exact interior is obtained from the disc and hole-complement interiors. A path's first boundary hit lies in the closure of its interior component (Lemmas~\ref{lem:fc-standard-inflation}--\ref{lem:fc-standard-component-boundary}). The actual continuous radial retraction is Theorem~\ref{thm:fc-radial-retraction}.

(3) Finite closed linear-coordinate boxes give neighbourhoods. Choose an upper disc of least radius, discard lower discs outside it, and retain maximal open holes; their nesting makes the retained holes disjoint. All inequalities at the original point are strict, so it is interior (Theorem~\ref{thm:fc-standard-affinoid-basis}). Whole-branch convergence to a Gauss point supplies a terminal annulus in every residue branch (Theorem~\ref{thm:fc-residue-end-annulus}).

(4) Algebraic roots and nontriviality derive dense valuation range. Actual evaluation points approach each positive-radius Gauss point from inside a branch; classification then puts every point in their closure (Theorem~\ref{thm:fc-rigid-density}). The whole-line contraction gives path connectedness and hence connectedness (Corollary~\ref{cor:fc-line-paths}).
\end{proof}

\begin{definition}[Functions on a standard affinoid]\label{def:fc-line-algebra}
\notready
\uses{def:fc-standard-affinoid}
$\mathcal O(U)$ is the completion of $F[T,(T-a_1)^{-1},\dots,(T-a_m)^{-1}]$ for $\|f\|_U=\max_{\zeta\in\partial U}|f|_\zeta$.
\end{definition}

\begin{lemma}[Functions]\label{lem:fc-line-functions}
\notready
\uses{def:fc-line-algebra, lem:fc-line-topology}
\begin{enumerate}
\item (Maximum modulus) For $f\in\mathcal O(U)$ and $x\in U$, $|f|_x\le\|f\|_U$. The points of $U$ are the bounded
multiplicative seminorms of $\mathcal O(U)$. Restriction maps exist for $U'\subseteq U$.
\item $f\in\mathcal O(U)$ is a unit if and only if $|f|_x\ne0$ for all $x\in U$.
\item (Gluing) If standard affinoids $U_i$ cover a closed disc $D$, then compatible $f_i\in\mathcal O(U_i)$ glue to some
$f\in\mathcal O(D)$.
\item (Liouville) If $c_R\in\mathcal O(D(0,R))$ are compatible for all $R$, and $\|c_R\|\le C\max(1,R)^N$, then they come from a
polynomial of degree $\le N$.
\end{enumerate}
\end{lemma}

\begin{proof}
(1) Mittag-Leffler decomposition $f=f_0(T)+\sum_jf_j((T-a_j)^{-1})$ with $\|f\|_U=\max(\|f_0\|_{\zeta_0},\|f_j\|_{\zeta_j})$, and
the ultrametric inequality at points of $U$. (2) Approximate $f$ by a rational $g$ with $|f-g|<\min_U|f|$. The zeros of
$g$ lie in holes or outside $D(a_0,r_0)$, so $1/g\in\mathcal O(U)$; then use a geometric series. (3) Induction over the
cover, the two-annulus case being \texttt{AnnulusGluing}. (4) Coefficient estimates $|c_n|R^n\le CR^N$.
\end{proof}

\begin{definition}[Harmonic functions]\label{def:fc-harmonic}
\notready
\uses{def:fc-standard-affinoid}
A function on $U$ is \emph{harmonic} if it has the form $c+\sum_{j=1}^m\gamma_j\log|T-a_j|$ with $c,\gamma_j\in\R$.
\end{definition}

\begin{lemma}[Harmonic functions]\label{lem:fc-harmonic}
\notready
\uses{def:fc-harmonic, lem:fc-line-topology}
Let $h$ be harmonic on $U$.
\begin{enumerate}
\item $h=h\circ\tau_U$, and $h$ is affine in $\log$-radius on edges of $\Sigma(U)$ and balanced at interior vertices.
\item $\max_Uh=\max_{\partial U}h$. At an interior local maximum, $h$ is constant on an actual neighbourhood. If an interior point realizes the whole-$U$ maximum, $h$ is constant on its connected component in $U\setminus\partial U$.
\item (Dirichlet) For all $(\beta_\zeta)\in\R^{\partial U}$ there is a unique harmonic $h$ with $h|_{\partial U}=\beta$.
\end{enumerate}
\end{lemma}

\begin{proof}
(1) Each term $\log|T-a_j|$ is $\log\max(r,|a-a_j|)$ at $\zeta_{a,r}$.
(2) The actual rational-unit local maximum theorem gives neighbourhood constancy.
For the weak maximum principle, the maximum locus is closed by continuity. If it
avoids the finite boundary, local constancy makes it open in the connected carrier,
a contradiction. The same open-and-closed argument gives constancy on the interior
component containing an interior maximum, as proved in the separate harmonic
boundary-comparison chapter. It does not imply whole-carrier constancy when hole
boundaries meet. (3) Index normalized logarithmic potentials by the actual distinct
nonouter Gauss boundary points, identifying coincident boundaries. Their independence
follows by choosing a nonzero coefficient with minimal hole radius and comparing
outward radial points before the next centre distance. Weak maximum applied to
$\pm h$ makes boundary evaluation injective. This is a square finite-dimensional
map on the actual boundary, so it is surjective; the outer constant supplies the
remaining boundary value. This is the actual Dirichlet construction, without the
incorrect dimension count $m+1$ when boundary points coincide.
\end{proof}

\begin{definition}[Subharmonic functions on the line {\cite[3.3]{GR25}}]\label{def:fc-line-subharmonic}
\lean{EnvelopeProperty.Berkovich.RelativeSpectrum.IsLineSubharmonicOn}
\leanok
\uses{def:fc-harmonic,def:fc-actual-line-subharmonic}
Let $\Omega\subseteq(\A^1)^\an$ be open. An usc $v\colon\Omega\to[-\infty,\infty)$ is \emph{subharmonic} if, for every standard
$U\subseteq\Omega$ and every harmonic $h$ on $U$, $v\le h$ on $\partial U$ implies $v\le h$ on $U$.
\end{definition}

\begin{lemma}[Basic properties]\label{lem:fc-line-sh-basic}
\lean{EnvelopeProperty.Berkovich.RelativeSpectrum.IsLineSubharmonicOn.add,EnvelopeProperty.Berkovich.RelativeSpectrum.IsLineSubharmonicOn.mul_nonneg,EnvelopeProperty.Berkovich.RelativeSpectrum.IsLineSubharmonicOn.sup,EnvelopeProperty.Berkovich.RelativeSpectrum.IsLineSubharmonicOn.finset_sup,EnvelopeProperty.Berkovich.RelativeSpectrum.IsLineSubharmonicOn.add_const,EnvelopeProperty.Berkovich.RelativeSpectrum.IsLineSubharmonicOn.add_harmonic,EnvelopeProperty.Berkovich.RelativeSpectrum.IsLineSubharmonicOn.mono,EnvelopeProperty.Berkovich.RelativeSpectrum.IsLineSubharmonicOn.iInf_directed,EnvelopeProperty.Berkovich.RelativeSpectrum.FCStandardAffinoid.isLineSubharmonicOn_analyticLog}
\leanok
\uses{def:fc-line-subharmonic,thm:fc-dirichlet-existence,thm:fc-completed-analytic-log-subharmonic}
\begin{enumerate}
\item Subharmonic functions are stable under sums, nonnegative multiples, finite maxima, addition of constants,
addition of harmonic functions, restriction to open subsets, and infima of decreasing nets indexed by a nonempty directed set.
\item $\log|f|$ is subharmonic on $\operatorname{int}(U')$ for every $f\in\mathcal O(U')$ in the actual completed affinoid ring.
\end{enumerate}
\end{lemma}

\begin{proof}
For sums, at every actual boundary point choose a real value $\beta$ between
$v_1$ and $h-v_2$; the exact extended-real splitting lemma permits bottom values.
The actual Dirichlet solution $H$ with these finite boundary values gives
$v_1\le H$ and $v_2\le h-H$ on the whole affinoid by comparison. Add the inequalities.
Nonnegative scaling follows by dividing a harmonic test for a positive scalar;
the zero scalar is handled directly. Maxima, constants and restrictions follow
from the literal comparison definition. Harmonic restriction and harmonic algebra
give addition of harmonic functions. For a nonempty decreasing directed infimum,
the finite boundary provides a common index for comparison against $h+\epsilon$;
then let $\epsilon$ decrease to zero. Upper semicontinuity of the infimum and
exclusion of positive infinity are proved without a lower bound.

For (2), original polynomial logarithms satisfy comparison by the actual rational
unit local-maximum theorem. A rational denominator has no zeros on the affinoid,
so its logarithm is actual harmonic. For an arbitrary completed function $f$ and
$\delta>0$, norm density constructs one original rational $g$ with
$\|f-g\|<\delta$. Contractive evaluation in every original completed residue field
and the ultrametric inequality give simultaneously on the whole affinoid
$\max(|f|,\delta)=\max(|g|,\delta)$. Its logarithm is subharmonic by the rational
case and maxima with constants. The directed infimum as $\delta$ decreases to zero
is exactly $\log|f|$, including negative infinity at zeros.
\end{proof}

\begin{lemma}[Strong maximum principle {\cite[Proposition~3.5]{GR25}}]\label{lem:fc-line-mp}
\lean{EnvelopeProperty.Berkovich.RelativeSpectrum.IsLineSubharmonicOn.eventually_eq_of_isLocalMax_norm}
\leanok
\uses{thm:fc-line-subharmonic-strong-maximum}
Under the standing nontrivially valued algebraically closed field convention,
if $v$ is subharmonic on the open $\Omega$ and has a local maximum $R$ at $x$,
then $v\equiv R$ on an actual neighbourhood of $x$.
\end{lemma}

\begin{proof}
The value $R=-\infty$ is immediate. Otherwise upper semicontinuity and actual
rigid density supply a lower rigid value if the function is not locally constant.
Construct a separated standard affinoid inside the local-maximum neighbourhood,
then append an actual low closed disc as a separated hole while retaining $x$
in the connected strict core. The actual Dirichlet solution with boundary values
$R$ and a strictly lower value on the new boundary is at most $R$ by weak maximum.
If it equals $R$ at $x$, local constancy makes it constant on the connected strict
core; actual outward Gauss approximation and continuity contradict its low boundary
value. Thus it is strictly lower at $x$, contradicting comparison. All separation,
low-hole preservation and strict-core closure properties are constructed internally
in the referenced checked strong-maximum theorem. A maximum relative to an open
$\Omega$ is an ambient local maximum because the neighbourhood filters agree there.
\end{proof}

\begin{lemma}[Locality]\label{lem:fc-line-local}
\lean{EnvelopeProperty.Berkovich.RelativeSpectrum.isLineSubharmonicOn_of_locally,EnvelopeProperty.Berkovich.RelativeSpectrum.isLineSubharmonicOn_iff_locally}
\leanok
\uses{thm:fc-line-subharmonic-locality}
Under the same standing valued-field convention, subharmonicity on an open domain
is local: if each point has an open neighbourhood on which the restriction is
subharmonic, the function is subharmonic on the whole domain.
\end{lemma}

\begin{proof}
Local upper semicontinuity and exclusion of positive infinity give the corresponding
global conditions. On any standard test affinoid, restrict its actual harmonic
function to each local neighbourhood, using the proved harmonic restriction theorem.
The local strong maximum principle gives local constancy at interior maxima of
the difference. The compact connected boundary-comparison argument therefore
propagates the boundary inequality across the whole test affinoid. The converse
is restriction of literal comparison to smaller open domains.
\end{proof}

\subsection{Finite flat curves over the line}\label{subsec:fc-pieces}

\begin{convention}\label{conv:fc-curve}
\notready
\uses{conv:fc-line}
$D$ is a smooth connected affine $F$-curve, and $p\colon D\to\A^1$ is finite flat of degree $d$, so $\mathcal O(D)$ is finite free over
$F[T]$. For a standard $W$, put $B_W=\mathcal O(W)\otimes_{F[T]}\mathcal O(D)$; its bounded seminorms are the points of $p^{-1}(W)$.
For $z$ over $w$, $m_z$ is the local multiplicity (\texttt{Berk.flatFibreMult}).
\end{convention}

\begin{lemma}[Pieces]\label{lem:fc-pieces}
\notready
\uses{conv:fc-curve, lem:fc-line-functions, lem:nf-pieces}
Let $w\in(\A^1)^\an$, $P\subseteq p^{-1}(w)$, and $N$ a neighbourhood of $P$ disjoint from $p^{-1}(w)\setminus P$. There are a
standard $W$ with $w\in\operatorname{int}W$ and an idempotent $e\in B_W$ such that $V\coloneqq\{|e|=1\}\subseteq N$ is clopen in
$p^{-1}(W)$ with $V\cap p^{-1}(w)=P$, and $eB_W$ is finite projective over $\mathcal O(W)$ of rank $m_P=\sum_{z\in P}m_z$. The local
norm $\Norm_V(b)=\det_{\mathcal O(W)}(b\mid eB_W)\in\mathcal O(W)$ satisfies $|\Norm_V(b)|_{w'}=\prod_{z\in V,\ p(z)=w'}|b|_z^{m_z}$ for all
$w'\in W$.
\end{lemma}

\begin{proof}
Approximate the idempotent of $P$ in $\mathcal H(w)\otimes\mathcal O(D)=\prod_zA_z$ by $b_0\in\mathcal O(D)$ (density). For $W$ small,
$\sigma(b_0^2-b_0)<1/4$ for the spectral seminorm $\sigma$ on $p^{-1}(W)$. Cayley--Hamilton over $\mathcal O(W)$ together with
the maximum modulus bounds the norms of powers by $\sigma$. Newton's iteration $b\mapsto3b^2-2b^3$ then converges in
$B_W$. The rank is constant because it equals $\sum_{z\in P}m_z$ at $w$ and is locally constant. The norm identity comes
from base change to $\mathcal H(w')$ (\texttt{SpectrumProjectiveNorm}).
\end{proof}

\begin{lemma}[Norms of convex PL functions]\label{lem:fc-local-norm-fs}
\notready
\uses{lem:fc-pieces, lem:fc-line-sh-basic, lem:finite-flat-norm}
In Lemma~\ref{lem:fc-pieces}, let $u=N^{-1}\max_j(\log|F_j|+\lambda_j)$, with $F_j$ regular and invertible on a Zariski
open set whose analytification contains $V$. Then
$\Norm_Vu(w')=\sum_{z\in V/w'}m_zu(z)=N^{-1}\max_\beta(\log|c_\beta|_{w'}+\langle\beta,\lambda\rangle)$, where the $c_\beta\in\mathcal O(W)$ are
the coefficients of $\Norm_V(\sum_jF_jT_j)$. In particular $\Norm_Vu$ is subharmonic on $\operatorname{int}W$.
\end{lemma}

\begin{proof}
Weighted Gauss norm identity, as in \cite[Lemma~5.13]{FHX}. The $F_j$ are units of $eB_W$ because their norms have no
zeros on $W$ (Lemma~\ref{lem:fc-line-functions}(2)). Then Lemma~\ref{lem:fc-line-sh-basic}.
\end{proof}

\subsection{Subharmonic and classically psh functions}\label{subsec:fc-gr25}

\begin{definition}[Subharmonic functions on curves {\cite[Definition~3.4]{GR25}}]\label{def:fc-subharmonic}
\notready
\uses{conv:fc-setting, lem:fc-pieces, def:fc-line-subharmonic}
Let $C$ be a curve over a complete field $F$, $U\subseteq C^\an$ open, and $u\colon U\to[-\infty,\infty)$ usc. A \emph{test datum} is
$(F',D,g,p,z,W,V)$, where:
\begin{itemize}
\item $F'/F$ is complete, algebraically closed and nontrivially valued;
\item $D$ is a smooth connected affine $F'$-curve and $g\colon D\to C_{F'}$ a morphism;
\item $p\colon D\to\A^1_{F'}$ is finite flat;
\item $z\in D^\an$ with $\pi g(z)\in U$;
\item $(W,V)$ is a piece (Lemma~\ref{lem:fc-pieces}) with $V\cap p^{-1}(p(z))=\{z\}$ and $\pi g(V)\subseteq U$.
\end{itemize}
$u$ is \emph{subharmonic} if, for every test datum, $\Norm_V(u\circ\pi\circ g)$ is subharmonic on $\operatorname{int}W$ or
$\equiv-\infty$ there.
\end{definition}

\begin{definition}[Classically psh functions {\cite[Definition~4.1]{GR25}}]\label{def:fc-classically-psh}
\notready
\uses{def:fc-subharmonic}
Let $X$ be an $F$-scheme of finite type and $U\subseteq X^\an$ open. An usc $u$ is \emph{classically psh} if $u\circ\pi\circ g$ is
subharmonic on $(\pi g)^{-1}(U)$ for every complete, algebraically closed, nontrivially valued $F'/F$ and every morphism
$g\colon D\to X_{F'}$ from a smooth connected affine $F'$-curve.
\end{definition}

\begin{proposition}[Properties of subharmonic functions {\cite[Proposition~3.5]{GR25}}]\label{prop:fc-subharmonic-properties}
\notready
\uses{def:fc-subharmonic, lem:fc-line-sh-basic, lem:fc-line-local, lem:fc-line-mp, lem:fc-pieces}
\begin{enumerate}
\item (Locality) Subharmonicity is local on $C^\an$.
\item Sums, nonnegative multiples and decreasing nets of subharmonic functions are subharmonic; so is $\log|f|$ for $f$
regular on a Zariski open set.
\item (Maximum principle) If $C$ is a smooth connected affine curve over a complete, algebraically closed, nontrivially
valued $F$, a subharmonic function with a local maximum $R$ at $z$ is $\equiv R$ near $z$.
\item If $u$ is subharmonic, so are $u\circ\pi_{F'/F}$ and $u\circ g'$ for every morphism of curves $g'$.
\item (Descent) If $\rho\colon C'\to C$ is finite surjective and $u\circ\rho$ is subharmonic, then $u$ is subharmonic.
\item (Norms, \cite[Lemma~5.21(1)]{FHX}) If $q\colon C'\to C$ is a finite flat morphism of normal curves and $u$ is subharmonic
on $C'^\an$, then $\Norm_qu$ is subharmonic on $C^\an$.
\end{enumerate}
\end{proposition}

\begin{proof}
(1) Refine pieces around each fibre point. $\Norm_V$ is then a sum of single-point norms, and
Lemma~\ref{lem:fc-line-local} applies. (2) Norms are additive and monotone, and finite weighted sums commute with
decreasing limits. $\Norm_V\log|f|=\log|\Norm_V(f)|$ (Lemma~\ref{lem:fc-line-sh-basic}(2)). (3) Take $p$ by Noether
normalization ($D$ smooth, so $p$ is flat), and a single-point piece $(W,V)$ at $z$ inside $\{u\le R\}$. Then
$\Norm_Vu\le m_zR$ with equality at $p(z)$, so $\Norm_Vu\equiv m_zR$ near $p(z)$ by Lemma~\ref{lem:fc-line-mp}. Since the
weights over each point sum to $m_z$ and every term is $\le R$, $u\equiv R$ on $V$ near $z$. (4) Test data compose.
(5) Let $D'$ be the normalization of a component of $D\times_CC'$ dominating $D$, of degree $\delta$ over $D$. Then the
local norms over $D$ equal $\delta^{-1}$ times the local norms over $D'$ (multiplicativity of multiplicities in
towers). (6) As in \cite[Lemma~5.21(1)]{FHX}, using flat base change of norms and transitivity of local norms.
\end{proof}

\begin{proposition}[Properties of classically psh functions {\cite[Proposition~4.2]{GR25}}]\label{prop:fc-cpsh-properties}
\notready
\uses{def:fc-classically-psh, prop:fc-subharmonic-properties}
Classically psh functions form a sheaf. They are stable under sums, nonnegative multiples, decreasing nets, forward
extension of scalars and pullback by morphisms of schemes of finite type. $\log|f|$ is classically psh for $f$ regular
on a Zariski open set. On curves, classically psh $=$ subharmonic.
\end{proposition}

\begin{proof}
Each property follows from the corresponding part of Proposition~\ref{prop:fc-subharmonic-properties}. On curves, the
two notions agree because subharmonicity is stable under pullback by morphisms of curves
(Proposition~\ref{prop:fc-subharmonic-properties}(4)).
\end{proof}

\begin{lemma}[Locally psh functions are classically psh]\label{lem:fc-locally-psh-cpsh}
\notready
\uses{def:fc-classically-psh, lem:fc-local-norm-fs, lem:local-basic, prop:fc-subharmonic-properties}
Every locally psh function on an open subset of $X^\an$ is classically psh. This is \cite[Lemma~5.2(4)]{FHX}.
\end{lemma}

\begin{proof}
Pull back to $D$ (Lemma~\ref{lem:local-basic}(2)). Near $z$ the pullback is a decreasing limit of convex PL functions
$N^{-1}\max_j(\log|F_{nj}|+\lambda_{nj})$, with $F_{nj}$ units on a Zariski open set; shrink the piece into that region.
Lemma~\ref{lem:fc-local-norm-fs} together with Lemma~\ref{lem:fc-line-sh-basic}(1) shows that $\Norm_V$ of the limit is
subharmonic.
\end{proof}

\begin{lemma}[Analytic curves are connected]\label{lem:fc-curve-connected}
\notready
\uses{lem:fc-pieces, lem:fc-line-functions, lem:fc-line-topology}
If $F$ is complete, algebraically closed and nontrivially valued, then $D^\an$ is connected for every irreducible affine
$F$-curve $D$.
\end{lemma}

\begin{proof}
Replace $D$ by its normalization, and take a separable finite flat $p\colon D\to\A^1$ with primitive $b$, so that $\chi_b$ is
irreducible over $F(T)$. Suppose $D^\an=A\sqcup B$ with $A,B$ clopen and nonempty. As $p$ is open and closed, $p(A)$ is all
of $(\A^1)^\an$, and likewise $p(B)$. For each $w$ take the piece with $P=A\cap p^{-1}(w)$ inside $A$: the characteristic
polynomial $\chi^A$ of $b$ on the pieces has coefficients in $\mathcal O(W)$. These coefficients glue on discs
(Lemma~\ref{lem:fc-line-functions}(3)) and have polynomial growth (integrality of $b$), so they are polynomials
(Lemma~\ref{lem:fc-line-functions}(4)). Then $\chi_b=\chi^A\chi^B$ is a nontrivial factorization in $F[T][Z]$, a
contradiction.
\end{proof}

\begin{proof}[Proof of Theorem~\ref{thm:fc-curve-mp} (Stage~2)]
\uses{lem:fc-locally-psh-cpsh, prop:fc-subharmonic-properties, lem:fc-curve-connected}
The function is classically psh (Lemma~\ref{lem:fc-locally-psh-cpsh}), hence subharmonic on $D^\an$. Its maximum set is
closed (usc) and open (Proposition~\ref{prop:fc-subharmonic-properties}(3)), and $D^\an$ is connected
(Lemma~\ref{lem:fc-curve-connected}).
\end{proof}

\begin{remark}[Comparison with \cite{GR25}]\label{rmk:fc-gr25}
\uses{def:fc-subharmonic, def:fc-classically-psh, thm:fc-global-mp}
\cite[Definition~3.4]{GR25} tests by Thuillier's subharmonicity on smooth strictly analytic curves; here the test is
by local norms to the line, and all test curves are algebraic. The two notions should agree, but a proof would use
Thuillier's theory on smooth curves, hence semistable reduction. It is not needed: every use in \cite{FHX} (Lemmas
5.2(4), 5.21 and Theorem~6.1, Step~1) passes through
Propositions~\ref{prop:fc-subharmonic-properties} and~\ref{prop:fc-cpsh-properties}. Not proved here:
\begin{itemize}
\item \cite[Prop.~3.5(3),(9)]{GR25};
\item the converse of 3.5(7);
\item stability under max of two arbitrary subharmonic functions;
\item \cite[Theorem~6.1]{GR25} at non-rigid points of higher-dimensional spaces.
\end{itemize}
Theorem~\ref{thm:fc-global-mp} is the consequence of the last item, together with \cite[3.4.8]{Berkovich90}, that
\cite{FHX} uses.
\end{remark}
